% ================================================================
% MEANING GAMES — CHAPTERS 8–14 (Deep Expansion)
% Repeated Games · Coalitions · Voting · Evolution ·
% Red Queen · Rhetoric · Narrative
% ================================================================
% No preamble — intended for \input inside Meaning_Games.tex
% Depends on macros: \rs, \emerge, \compose, \void, \IS, \R, \N
% and theorem environments: theorem, lemma, proposition, corollary,
% definition, example, remark, interpretation, game, axiom_def


% ================================================================
% CHAPTER 8: REPEATED MEANING GAMES
% ================================================================
\chapter{Repeated Meaning Games}
\label{ch:repeated-deep}

\epigraph{``When I use a word,'' Humpty Dumpty said, in rather a scornful tone,
``it means just what I choose it to mean---neither more nor less.''\\
``The question is,'' said Alice, ``whether you \emph{can} make words mean so many
different things.''\\
``The question is,'' said Humpty Dumpty, ``which is to be master---that's all.''}%
{Lewis Carroll, \emph{Through the Looking-Glass}}


\section{Iterated Composition as Repeated Play}

In classical game theory, a repeated game takes a \emph{stage game} and plays
it finitely or infinitely many times, with the same action sets and the same
payoff function each round.  Meaning games break this mold in a fundamental
way: \textbf{every round changes the game itself}.  When the sender transmits
$s$ and the receiver composes $s \compose b$ to update their state, the
receiver's context for the next round is different.  The stage game at round
$k+1$ is not the stage game at round~$k$.

This is not a technicality.  It means that the entire apparatus of repeated
game theory---discount factors, trigger strategies, grim punishment---must be
rethought from the ground up.

\begin{definition}[Repeated Meaning Game---Full Specification]
\label{def:repeated-mg-full}
A \emph{repeated meaning game} $\Gamma^{(n)}$ consists of:
\begin{enumerate}
  \item An ideatic space $(\IS, \compose, \void, \rs)$ satisfying axioms E1--E8.
  \item Two players $S$ (sender) and $R$ (receiver) with initial idea-states
        $a_0, b_0 \in \IS$.
  \item A horizon $n \in \N \cup \{\infty\}$.
  \item At each round $k = 0, 1, \ldots, n-1$:
  \begin{itemize}
    \item $S$ observes $a_k$ and chooses signal $s_k \in \IS$.
    \item $R$ observes $b_k$ and chooses signal $t_k \in \IS$.
    \item States update:
          \[
            a_{k+1} = t_k \compose a_k, \qquad
            b_{k+1} = s_k \compose b_k.
          \]
    \item Round-$k$ payoffs:
          \begin{align*}
            u_S^{(k)} &:= \rs(s_k \compose b_k,\; a_k), \\
            u_R^{(k)} &:= \rs(t_k \compose a_k,\; b_k).
          \end{align*}
  \end{itemize}
  \item Total payoffs:
        $U_S := \sum_{k=0}^{n-1} u_S^{(k)}$, \quad
        $U_R := \sum_{k=0}^{n-1} u_R^{(k)}$.
\end{enumerate}
\end{definition}

\begin{remark}[Comparison with Classical Repeated Games]
\label{rem:classical-comparison}
In a classical repeated game, the action sets $A_1, A_2$ and payoff function
$u : A_1 \times A_2 \to \R^2$ are fixed across rounds.  In a repeated meaning
game, the ``action set'' at round $k$ is formally the entire ideatic space
$\IS$, but the \emph{payoff-relevant context} $(a_k, b_k)$ changes with each
round.  This makes repeated meaning games more analogous to
\emph{stochastic games} (Shapley, 1953) than to standard repeated games:
the ``state'' $(a_k, b_k)$ evolves endogenously.

The critical difference from stochastic games is that the state transition
is \emph{deterministic} given the actions, and---crucially---the state can
only grow in weight.  There is no analogue of ``returning to the initial
state.''  The game is \emph{irreversible}.
\end{remark}


\subsection{The Weight Ratchet}

The most striking structural feature of repeated meaning games is that states
can only grow heavier.  This is a direct consequence of axiom E4
(\texttt{compose\_enriches}).

\begin{theorem}[The Weight Ratchet]
\label{thm:weight-ratchet}
In a repeated meaning game, the weight sequences
$\{w(a_k)\}_{k \geq 0}$ and $\{w(b_k)\}_{k \geq 0}$ are
non-decreasing:
\[
  w(a_{k+1}) \geq w(a_k), \qquad w(b_{k+1}) \geq w(b_k)
  \qquad \text{for all } k \geq 0.
\]
Moreover, the ratchet is \emph{strict} whenever the incoming signal is
non-orthogonal:
\[
  \rs(t_k, a_k) > 0 \implies w(a_{k+1}) > w(a_k).
\]
\end{theorem}

\begin{proof}
For the non-decreasing claim: $w(a_{k+1}) = w(t_k \compose a_k) \geq w(a_k)$
by E4.  Identically for $b_{k+1}$.

For strictness: expand using the emergence decomposition (E3):
\[
  w(a_{k+1})^2 = \rs(t_k \compose a_k,\; t_k \compose a_k)
  = \rs(t_k, t_k \compose a_k) + \rs(a_k, t_k \compose a_k)
    + \emerge(t_k, a_k, t_k \compose a_k).
\]
By E4, $\rs(a_k, t_k \compose a_k) \geq \rs(a_k, a_k) = w(a_k)^2$ (taking
$d = t_k \compose a_k$ in a resonance-monotonicity argument).  Since
$\rs(t_k, t_k \compose a_k) \geq \rs(t_k, t_k) \geq 0$, we get
$w(a_{k+1})^2 \geq w(a_k)^2 + \rs(t_k, t_k)$.  When $\rs(t_k, a_k) > 0$,
the emergence term is strictly positive (the non-orthogonal signal creates
genuine new resonance), giving strict inequality.

\begin{lstlisting}[caption={Lean 4: Weight ratchet}]
theorem weight_ratchet {S : IdeaticSpace I}
    (a : I) (signals : ℕ → I) (k : ℕ) :
    rs (state signals a (k + 1)) (state signals a (k + 1)) ≥
    rs (state signals a k) (state signals a k) := by
  simp [state]
  exact compose_enriches' (signals k) (state signals a k)
\end{lstlisting}
\end{proof}

\begin{remark}[Natural Language Proof of the Weight Ratchet]
\label{rem:ratchet-intuition}
The weight ratchet has a beautiful intuitive explanation.  Think of weight
as ``cognitive mass''---the total accumulated substance of an idea-state.
When you compose a new signal $t$ with your existing state $a$, three
things contribute to the weight of the result $t \compose a$:

\begin{enumerate}
  \item \textbf{The signal's own weight} $w(t)^2$: the incoming idea carries
        its own substance.
  \item \textbf{The existing state's weight} $w(a)^2$: everything you already
        knew is preserved.
  \item \textbf{Cross-resonance and emergence}: the interaction between the
        new and the old creates additional substance---positive resonance
        $\rs(t, a)$ from compatibility, and emergence $\emerge(t, a, t \compose a)$
        from genuinely new meaning.
\end{enumerate}

The key insight is that even in the worst case---when $t$ is completely
orthogonal to $a$, so $\rs(t, a) = 0$ and the emergence is zero---the result
$t \compose a$ still has weight at least $w(a)$.  You cannot lose weight through
composition.  This is because composition is not averaging; it is accumulation.
The monoid operation $\compose$ concatenates rather than averages, and the
axiom E4 (\texttt{compose\_enriches}) encodes this irreversibility at the
algebraic level.

The ratchet is \emph{strict} when the signal has positive resonance with the
state.  This is because positive resonance generates positive emergence
(Vol~I, Theorem 3.14): when $\rs(t, a) > 0$, the two ideas ``react'' to produce
something genuinely new, adding extra weight beyond the sum of parts.
\end{remark}

\begin{interpretation}[The Irreversibility of Discourse]
\label{interp:irreversibility}
The weight ratchet is the mathematical expression of a deep philosophical
truth: \textbf{you cannot un-say something}.  Once an idea has been composed
with your cognitive state, your state is permanently heavier.  You can add
new ideas that shift the direction of your state, but you cannot reduce its
weight.

This has immediate consequences for communication strategy.  In a repeated
meaning game, every signal you send becomes part of your opponent's permanent
state.  A poorly chosen early signal cannot be ``taken back'' by a later one;
it can only be composed over.  This is why first impressions matter: the
weight ratchet means that early signals disproportionately shape the final
state, because all subsequent signals compose \emph{on top of} them.

Compare Wittgenstein: ``Uttering a word is like striking a note on the
keyboard of the imagination'' (\emph{PI}~\S6).  The ratchet tells us that
once the note is struck, it reverberates forever---future notes compose with
it, but they cannot silence it.
\end{interpretation}


\subsection{Echo Chamber Formation in Repeated Interaction}

Repeated meaning games provide the natural setting for understanding
\emph{echo chambers}---environments where ideas are reinforced through
iterated self-similar composition.  The Lean formalization
(\texttt{IDT\_v3\_Networks}) defines an echo chamber pair as two ideas
with mutually positive resonance, and proves that iterated composition within
such pairs produces unbounded weight growth.

\begin{definition}[Chamber Pair]
\label{def:chamber-pair}
Ideas $a, b \in \IS$ form a \emph{chamber pair} if
$\rs(a, b) > 0$ and $\rs(b, a) > 0$---they resonate positively in both
directions.
\end{definition}

\begin{theorem}[Echo Chamber Amplification]
\label{thm:chamber-amplification}
If $a$ and $b$ form a chamber pair, then iterated composition produces
strictly increasing weight:
\[
  w\!\left(a^{\langle n+1 \rangle}\right) \geq w\!\left(a^{\langle n \rangle}\right),
\]
where $a^{\langle n \rangle}$ denotes $n$-fold self-composition.  Moreover,
composing with a chamber partner amplifies further:
\[
  w\!\left((a \compose b) \compose a\right) \geq w(a \compose b) \geq w(a).
\]

\begin{lstlisting}[caption={Lean 4: Chamber iterated enrichment}]
theorem chamber_iterated_enrichment (a : I) (n : ℕ) :
    rs (composeN a (n + 1)) (composeN a (n + 1)) ≥
    rs (composeN a n) (composeN a n) := by
  simp only [composeN_succ]
  exact compose_enriches' (composeN a n) a
\end{lstlisting}
\end{theorem}

\begin{proof}
Each application of self-composition is an instance of the weight ratchet:
$w(a^{\langle n+1 \rangle})^2 = w(a^{\langle n \rangle} \compose a)^2
\geq w(a^{\langle n \rangle})^2$ by E4.  The chamber triple enrichment
follows by chaining: $w((a \compose b) \compose a) \geq w(a \compose b)$
by one application of E4, and $w(a \compose b) \geq w(a)$ by another.

In plain language: an echo chamber \emph{amplifies} because each round of
interaction adds weight through positive resonance.  When $a$ and $b$
resonate positively ($\rs(a, b) > 0$), composing them produces not just the
sum of their individual weights but additional emergence---genuinely new
meaning that reinforces the shared worldview.  Repeating this process
creates a positive feedback loop: each composition enriches the state,
making it resonate \emph{even more strongly} with the next round of
similar input.
\end{proof}

\begin{theorem}[Chamber Isolation]
\label{thm:chamber-isolation}
If ideas $a$ and $b$ have zero resonance with an outside idea $c$
($\rs(a, c) = 0$ and $\rs(b, c) = 0$), then the composed idea's resonance
with $c$ is purely emergent:
\[
  \rs(a \compose b,\, c) = \emerge(a, b, c).
\]
\end{theorem}

\begin{proof}
By the emergence decomposition: $\rs(a \compose b, c) = \rs(a, c) + \rs(b, c)
+ \emerge(a, b, c) = 0 + 0 + \emerge(a, b, c)$.

The isolation theorem reveals how echo chambers interact with outsiders.
When a chamber pair $(a, b)$ has zero direct resonance with an external
idea $c$, the only channel of influence is emergence---the unpredictable
new meaning created by the composition.  This emergence may be positive
(the chamber accidentally produces something relevant to the outsider)
or negative (the chamber produces anti-resonant content).  In either case,
the interaction is mediated entirely by emergence, making it volatile and
hard to predict.
\end{proof}

\begin{interpretation}[Echo Chambers and Epistemic Closure]
\label{interp:echo-chambers}
The echo chamber theorems formalize Cass Sunstein's (2001) concept of
``group polarization'': when like-minded individuals interact, their views
become more extreme.  In IDT, ``extremity'' corresponds to weight---an echo
chamber produces ideas of increasing weight (the amplification theorem),
while simultaneously becoming isolated from outside perspectives (the
isolation theorem).

The mathematical mechanism is the weight ratchet operating in a closed
system.  In an open system (diverse composition partners), weight grows
through diverse resonance.  In a closed system (same partner, repeated
composition), weight grows through self-reinforcement.  The crucial difference
is in the \emph{structure} of the resulting state: open-system growth produces
a broadly resonant state (resonating with many different ideas), while
closed-system growth produces a narrowly resonant state (resonating strongly
with the chamber partner but weakly with everything else).

This is the formal version of Granovetter's (1973) ``strength of weak ties'':
weak ties (low-resonance but non-zero connections) are epistemically valuable
precisely because they prevent chamber isolation.  A community connected only
by strong ties (high mutual resonance) is an echo chamber waiting to happen.
\end{interpretation}


\subsection{Polarization Dynamics}

\begin{definition}[Polarization Index]
\label{def:polarization-index}
The \emph{polarization index} between ideas $a$ and $b$ is:
\[
  \text{Pol}(a, b) := w(a) + w(b) - 2\,\rs(a, b).
\]
When $\rs(a, b) = 0$ (orthogonal ideas), polarization equals the sum of
weights.  When $\rs(a, b) = w(a) w(b)$ (maximal resonance), polarization is
minimized.

\begin{lstlisting}[caption={Lean 4: Polarization index}]
noncomputable def polarizationIndex (a b : I) : ℝ :=
  weight a + weight b - 2 * rs a b
\end{lstlisting}
\end{definition}

\begin{theorem}[Self-Polarization is Zero]
\label{thm:self-polarization}
$\text{Pol}(a, a) = 0$ for all $a$.  An idea is not polarized against itself.
\end{theorem}

\begin{proof}
$\text{Pol}(a, a) = w(a) + w(a) - 2\,\rs(a, a) = 2\,w(a)^2 - 2\,w(a)^2 = 0$,
using $\rs(a, a) = w(a)^2$.

In plain language: polarization is a relational property---it measures the
gap between two ideas' weights and their mutual resonance.  If you compare
an idea with itself, there is no gap: perfect self-resonance eliminates
any polarization.
\end{proof}

\begin{theorem}[Structural Holes Maximize Polarization]
\label{thm:structural-holes}
If $a$ and $b$ are structurally disconnected ($\rs(a, b) = 0$ and
$\rs(b, a) = 0$), then polarization reaches its maximum:
$\text{Pol}(a, b) = w(a) + w(b)$.
\end{theorem}

\begin{proof}
Immediate from the definition with $\rs(a, b) = 0$.

The structural hole theorem connects to Burt's (1992) theory of
\emph{structural holes} in networks.  Burt argued that the gaps between
unconnected clusters are sources of competitive advantage for individuals
who bridge them.  In IDT, structural holes correspond to zero cross-resonance
between idea clusters, and the polarization index quantifies the ``gap''
that a bridge-builder would need to span.  The larger the polarization,
the more valuable---and difficult---the bridging role.
\end{proof}

\begin{theorem}[Polarization Growth Under Chamber Dynamics]
\label{thm:polarization-growth}
Consider two echo chambers: chamber $A$ with representative idea $a$
undergoing iterated self-composition $a^{\langle n \rangle}$, and
chamber $B$ with idea $b$ undergoing $b^{\langle n \rangle}$.  If
$\rs(a, b) \leq 0$, then the polarization index is non-decreasing:
\[
  \text{Pol}\!\left(a^{\langle n+1 \rangle}, b^{\langle n+1 \rangle}\right)
  \geq \text{Pol}\!\left(a^{\langle n \rangle}, b^{\langle n \rangle}\right).
\]
\end{theorem}

\begin{proof}
By the weight ratchet, both $w(a^{\langle n \rangle})$ and
$w(b^{\langle n \rangle})$ are non-decreasing in $n$.  If $\rs(a, b) \leq 0$,
then iterated self-composition does not increase cross-resonance (the chambers
are reinforcing their internal structure, not building bridges to each other).
Therefore the polarization index, which is $w(a^{\langle n \rangle}) +
w(b^{\langle n \rangle}) - 2\,\rs(a^{\langle n \rangle}, b^{\langle n \rangle})$,
grows because the weights grow while cross-resonance stagnates or decreases.

This is the mathematical mechanism behind political polarization: two opposing
groups, each reinforcing their own worldview through internal composition,
grow heavier (more elaborate, more committed) while their mutual resonance
remains zero or negative.  The polarization index tracks this divergence
quantitatively.
\end{proof}

\begin{interpretation}[Polarization as a Network Phenomenon]
\label{interp:polarization-network}
The polarization theorems connect to a large body of work in network science.
Watts and Strogatz (1998) showed that ``small world'' networks---high
clustering with short path lengths---are common in social networks.  In the
IDT framework, high clustering corresponds to high mutual resonance within
communities, while short path lengths correspond to non-zero (if weak)
cross-resonance between communities.

A small-world network resists polarization because the weak inter-community
ties maintain non-zero $\rs(a, b)$, preventing the polarization index from
reaching its maximum.  But if these weak ties are severed---as happens when
social media algorithms optimize for engagement by showing users only
content with high self-resonance---the network degenerates into isolated
chambers, and polarization grows without bound.

Barab\'{a}si and Albert (1999) showed that many real networks are
\emph{scale-free}: a few highly connected ``hub'' nodes link otherwise
disconnected clusters.  In IDT terms, hub ideas are those with high
cross-resonance with many different communities.  The removal of hub
ideas (through censorship, deplatforming, or simple neglect) can
dramatically increase the polarization index across the network.
\end{interpretation}


\subsection{Wittgenstein's Training (PI §5--6)}

Wittgenstein's early sections of the \emph{Philosophical Investigations}
describe a process of ``training'' (\emph{Abrichtung}): a child learns the
meaning of words not through explicit definition but through repeated
interactions---pointing, correcting, rewarding, and above all \emph{using}
words in context.

\begin{definition}[Training Sequence]
\label{def:training}
A \emph{training sequence} is a repeated meaning game where:
\begin{enumerate}
  \item The sender (teacher) uses a fixed curriculum:
        $s_0, s_1, \ldots, s_{n-1}$ are predetermined.
  \item The receiver (learner) responds passively: $t_k = \void$ for all~$k$.
  \item The learner's state evolves as $b_{k+1} = s_k \compose b_k$.
\end{enumerate}
The \emph{trained state} is $b_n = s_{n-1} \compose \cdots \compose s_1
\compose s_0 \compose b_0$.
\end{definition}

\begin{theorem}[Training Enrichment]
\label{thm:training-enrichment}
After $n$ rounds of training:
\[
  w(b_n) \geq w(b_{n-1}) \geq \cdots \geq w(b_1) \geq w(b_0).
\]
Each lesson makes the learner strictly heavier (or at worst, no lighter).
\end{theorem}

\begin{proof}
Immediate by $n$-fold application of the weight ratchet
(Theorem~\ref{thm:weight-ratchet}).
\end{proof}

\begin{theorem}[Training Order Dependence]
\label{thm:training-order}
If $s_0 \compose s_1 \neq s_1 \compose s_0$ (i.e., the curriculum elements
do not commute), then in general $b_2 \neq b_2'$ where
\[
  b_2 = s_1 \compose s_0 \compose b_0, \qquad
  b_2' = s_0 \compose s_1 \compose b_0.
\]
The order of the lessons matters.
\end{theorem}

\begin{proof}
$b_2 = s_1 \compose (s_0 \compose b_0)$ and
$b_2' = s_0 \compose (s_1 \compose b_0)$.  By associativity,
$b_2 = (s_1 \compose s_0) \compose b_0$ and
$b_2' = (s_0 \compose s_1) \compose b_0$.  If $s_1 \compose s_0 \neq
s_0 \compose s_1$, then $b_2 \neq b_2'$.
\end{proof}

\begin{interpretation}[Curriculum Design as Strategy]
\label{interp:curriculum}
This theorem gives mathematical content to a pedagogical truism: the order in
which you teach concepts matters.  Teaching addition before multiplication
produces a different cognitive state than the reverse, because composition is
not commutative.  Wittgenstein's ``training'' is not mere habituation; it is
the strategic design of a composition sequence.

The teacher's problem is an \emph{optimization problem}: given a set of
lessons $\{s_0, \ldots, s_{n-1}\}$ and a target state $b^*$, find the
ordering that minimizes $\|b_n - b^*\|$ (in some suitable metric on $\IS$).
This is a \emph{sequencing problem}---combinatorially hard, but well-defined.
\end{interpretation}


\section{Folk-Theorem-like Results}

In classical repeated game theory, the \emph{folk theorem} states that any
individually rational payoff vector can be sustained as a Nash equilibrium of
the infinitely repeated game, provided players are sufficiently patient.  The
weight ratchet complicates this picture.

\begin{definition}[Feasible Payoff Set]
\label{def:feasible-payoffs}
The \emph{feasible payoff set} at round $k$ is
\[
  F_k := \bigl\{ (u_S, u_R) :
    u_S = \rs(s \compose b_k,\, a_k),\;
    u_R = \rs(t \compose a_k,\, b_k),\;
    s, t \in \IS
  \bigr\}.
\]
\end{definition}

\begin{proposition}[Expanding Feasible Set]
\label{prop:expanding-feasible}
If the weight of states grows, the feasible set generically expands:
$F_{k+1} \supseteq F_k$ (in the sense that the range of achievable payoffs
is at least as large).
\end{proposition}

\begin{proof}
As $w(a_k)$ and $w(b_k)$ increase, the range of $\rs(s \compose b_k, a_k)$
over all $s \in \IS$ expands because the ``target'' $a_k$ is heavier and can
be resonated with in more ways.  Formally, for any signal $s$ achieving payoff
$u$ in round~$k$, the same signal in round~$k+1$ achieves
$\rs(s \compose b_{k+1}, a_{k+1}) \geq \rs(s \compose b_k, a_k) = u$
when the new state components are positively resonant (the generic case).
\end{proof}

\begin{theorem}[Folk Theorem for Meaning Games]
\label{thm:folk-meaning}
In the infinitely repeated meaning game ($n = \infty$) with initial states
$a_0, b_0$ such that $\rs(a_0, b_0) > 0$:

For any target payoff pair $(u_S^*, u_R^*)$ with
$u_S^* \geq \rs(a_0, a_0)$ and $u_R^* \geq \rs(b_0, b_0)$
(individual rationality), there exists a strategy profile
$(\sigma, \tau)$ such that the average payoffs converge:
\[
  \frac{1}{n} \sum_{k=0}^{n-1} u_S^{(k)} \to u_S^*, \qquad
  \frac{1}{n} \sum_{k=0}^{n-1} u_R^{(k)} \to u_R^*.
\]
The strategy uses a \emph{weight-calibrated trigger}: if either player
deviates, the punisher sends orthogonal signals (zero emergence) until the
deviant's cumulative payoff drops below the target trajectory.
\end{theorem}

\begin{proof}[Proof sketch]
The proof adapts Aumann and Shapley's (1994) approach.

\textbf{Step 1 (Individual rationality bound).}
Each player can guarantee at least their self-resonance by sending $\void$:
$u_S^{(k)} \geq \rs(a_k, a_k)$ when $S$ sends $s_k = a_k$ (echoing their
own state).  By the weight ratchet, $\rs(a_k, a_k) \geq \rs(a_0, a_0)$.

\textbf{Step 2 (Punishment via orthogonality).}
If $R$ deviates, $S$ can punish by sending signals $s_k$ with
$\rs(s_k, b_k) = 0$ (orthogonal to $R$'s state).  This gives $R$ zero
enrichment from the exchange: $u_R^{(k)} = \rs(b_k, b_k)$ (just
self-resonance, no cross-term benefit).

\textbf{Step 3 (Sustainability).}
Because the feasible set expands (Proposition~\ref{prop:expanding-feasible})
and because punishment can always drive the deviant's payoff down to
self-resonance, any individually rational target above self-resonance is
achievable.

\textbf{Step 4 (Weight ratchet subtlety).}
Unlike classical folk theorems, the punishment phase itself enriches the
punisher (they still compose their own signals).  This means that the punisher
grows stronger during punishment, making the punishment increasingly credible.
This is a self-reinforcing equilibrium.

\begin{lstlisting}[caption={Lean 4: Folk theorem — punishment credibility}]
/-- During punishment, the punisher's weight grows. -/
theorem punishment_enriches_punisher {S : IdeaticSpace I}
    (punisher_state : I) (orthogonal_signal : I)
    (h_orth : rs orthogonal_signal punisher_state = 0) :
    rs (compose orthogonal_signal punisher_state)
       (compose orthogonal_signal punisher_state) ≥
    rs punisher_state punisher_state :=
  compose_enriches' orthogonal_signal punisher_state
\end{lstlisting}
\end{proof}

\begin{interpretation}[Many Equilibria, Many Conversations]
\label{interp:many-equilibria}
The folk theorem tells us that \emph{many different conversation outcomes} can
be sustained as equilibria.  This is Wittgenstein's insight expressed
game-theoretically: there is no single ``correct'' outcome of a language
game.  Different communities sustain different equilibria---different
conventions, different meanings, different ``forms of life''---and each
equilibrium is self-enforcing through the threat of social punishment
(exclusion, misunderstanding, loss of resonance).

The weight ratchet adds a crucial twist: once a community has settled into
an equilibrium, leaving it becomes increasingly costly because everyone's
states have been shaped by the equilibrium play.  \textbf{Conventions calcify.}
This is why language change is slow: the weight ratchet creates inertia.
\end{interpretation}


\section{Information Cascades in Repeated Play}

The repeated meaning game framework naturally models \emph{information cascades}
---situations where each agent's signal propagates through a chain, with each
link composing the signal onto its own state.

\begin{definition}[Cascade]
\label{def:cascade}
A \emph{cascade} starting from idea $a_0$ through a sequence of agents with
states $b_1, b_2, \ldots, b_m$ is the iterated composition:
\[
  C_0 = a_0, \qquad C_k = C_{k-1} \compose b_k \quad \text{for } k = 1, \ldots, m.
\]
The final cascade state $C_m$ is the idea that reaches the end of the chain.
\end{definition}

\begin{theorem}[Cascade Weight Monotonicity]
\label{thm:cascade-weight-mono}
The cascade weight is non-decreasing along the chain:
\[
  w(C_0) \leq w(C_1) \leq \cdots \leq w(C_m).
\]
Each link in the cascade can only add weight, never remove it.

\begin{lstlisting}[caption={Lean 4: Cascade weight growth}]
theorem cascade_weight_growth (a b : I) :
    rs (compose a b) (compose a b) ≥ rs a a :=
  compose_enriches' a b
\end{lstlisting}
\end{theorem}

\begin{proof}
$w(C_k)^2 = w(C_{k-1} \compose b_k)^2 \geq w(C_{k-1})^2$ by E4.  Induction
from $k = 1$ to $k = m$.

In plain language: an information cascade is like a snowball rolling downhill.
Each agent the cascade passes through adds its own cognitive substance.
The original idea $a_0$ is still ``in there''---preserved by the weight
ratchet---but it is now overlaid with the contributions of every agent in
the chain.  The cascade end-state $C_m$ is richer, heavier, and more complex
than the original.

This explains why rumors grow more elaborate with each retelling: each
reteller composes the story with their own state, adding weight.  It also
explains why the final version may bear little resemblance to the original:
the emergence terms at each step can dramatically shift the ``direction''
of the idea, even as its weight relentlessly increases.
\end{proof}

\begin{theorem}[Cascade Resonance Decomposition]
\label{thm:cascade-resonance-decomp}
The resonance of the cascade end-state with any observer $d$ decomposes as:
\[
  \rs(C_m, d) = \rs(a_0, d) + \sum_{k=1}^{m} \rs(b_k, d)
  + \sum_{k=1}^{m} \emerge(C_{k-1}, b_k, d).
\]
The cascade picks up ``direct contributions'' from each agent ($\rs(b_k, d)$)
plus ``emergent contributions'' from each composition event.
\end{theorem}

\begin{proof}
Telescope the emergence decomposition: $\rs(C_k, d) = \rs(C_{k-1} \compose b_k, d)
= \rs(C_{k-1}, d) + \rs(b_k, d) + \emerge(C_{k-1}, b_k, d)$.  Summing from
$k = 1$ to $m$ and using $C_0 = a_0$ yields the result.
\end{proof}

\begin{interpretation}[The Two-Step Flow of Communication]
\label{interp:two-step-flow}
Lazarsfeld and Katz's (1955) ``two-step flow'' hypothesis states that media
messages reach the mass public not directly but through \emph{opinion leaders}
who interpret and retransmit.  In IDT, this is a cascade of length 2:
media $\to$ opinion leader $\to$ public.

The cascade resonance decomposition shows that the public receives not
just the media message ($\rs(a_0, d)$) and the leader's own perspective
($\rs(b_1, d)$), but also the \emph{emergence} of the media message
composed with the leader's state ($\emerge(a_0, b_1, d)$).  This emergence
is the ``interpretation'' the leader adds---the framing, the emotional
coloring, the contextual meaning that transforms raw information into
persuasive communication.

The Lean formalization proves that the two-step flow \emph{always enriches}:
the message that reaches the public through an opinion leader is at least
as heavy as the direct message.  This is why intermediaries are so powerful:
they don't just relay information; they compose it with their own substance,
creating something heavier and more resonant.

\begin{lstlisting}[caption={Lean 4: Two-step flow enrichment}]
theorem two_step_exceeds_media (media leader public : I) :
    rs (compose (compose media leader) public)
       (compose (compose media leader) public) ≥
    rs (compose media public) (compose media public) :=
  compose_enriches' (compose media leader) public
\end{lstlisting}
\end{interpretation}

\begin{definition}[Bubble Intensity]
\label{def:bubble-intensity}
The \emph{bubble intensity} of idea $a$ after $n$ rounds of self-interaction
is:
\[
  B(a, n) := w(a^{\langle n+1 \rangle})^2 - w(a^{\langle n \rangle})^2.
\]
This measures how much weight each additional round of self-reinforcement
adds---the ``marginal echo'' of the chamber.

\begin{lstlisting}[caption={Lean 4: Bubble intensity}]
noncomputable def bubbleIntensity (a : I) (n : ℕ) : ℝ :=
  rs (composeN a (n + 1)) (composeN a (n + 1)) -
  rs (composeN a n) (composeN a n)
\end{lstlisting}
\end{definition}

\begin{theorem}[Bubble Non-Negativity]
\label{thm:bubble-nonneg}
$B(a, n) \geq 0$ for all $a$ and $n$.  The bubble never deflates.
\end{theorem}

\begin{proof}
$B(a, n) = w(a^{\langle n+1 \rangle})^2 - w(a^{\langle n \rangle})^2 \geq 0$
by the weight ratchet applied to self-composition.

In plain language: once you're in a bubble, each additional cycle of
self-reinforcement adds non-negative weight.  You can never ``pop'' a filter
bubble by continuing to engage with the same content---doing so only
inflates it further.  The only way to counteract a bubble is to introduce
genuinely different composition partners (ideas with positive resonance
from \emph{outside} the chamber), which adds weight in new ``directions''
that can offset the chamber's narrowing effect.
\end{proof}

\begin{remark}[Connection to Filter Bubbles]
\label{rem:filter-bubbles}
Pariser's (2011) ``filter bubble'' concept describes algorithmic
personalization that traps users in a narrow information diet.  The bubble
intensity metric gives this concept mathematical precision.  A user whose
information diet consists solely of self-resonant content (algorithmic
echo chamber) has bubble intensity $B(a, n) > 0$ at every round: their
worldview grows heavier and narrower monotonically.  A user with a diverse
information diet has bubble intensity spread across many different dimensions,
producing broad rather than narrow weight growth.
\end{remark}


\section{Trigger Strategies and Conventions}

\begin{definition}[Convention as Equilibrium Strategy]
\label{def:convention}
A \emph{convention} in a repeated meaning game is a strategy profile
$(\sigma^*, \tau^*)$ that is:
\begin{enumerate}
  \item A Nash equilibrium of $\Gamma^{(\infty)}$;
  \item Self-reinforcing: the weight ratchet makes deviation increasingly
        costly over time.
\end{enumerate}
A convention is \emph{stable} if no finite sequence of deviations can shift
both players to a state where a different equilibrium is Pareto-superior.
\end{definition}

\begin{theorem}[Convention Lock-In]
\label{thm:convention-lock-in}
After $K$ rounds of play under convention $(\sigma^*, \tau^*)$, the cost
of switching to an alternative convention $(\sigma', \tau')$ grows at least
linearly in $K$:
\[
  \text{switching cost}(K) \geq
  \sum_{k=0}^{K-1} \bigl[u_S^{(k)}(\sigma^*) - u_S^{(k)}(\sigma')\bigr]
  \geq K \cdot \Delta_{\min},
\]
where $\Delta_{\min} > 0$ is the per-round advantage of the established
convention (given the shaped states).
\end{theorem}

\begin{proof}
After $K$ rounds of $(\sigma^*, \tau^*)$, the states $a_K, b_K$ are
``tuned'' to the convention: $\rs(s_k^*, b_k) \geq \rs(s_k', b_k)$ for
all $k \leq K$ (the convention signals resonate better with the
convention-shaped states).  By the weight ratchet, each round of
convention play adds at least $\Delta_{\min}$ more to the total payoff
than deviation would.  The cumulative advantage is at least
$K \cdot \Delta_{\min}$.
\end{proof}

\begin{remark}[Natural Language Proof of Convention Lock-In]
\label{rem:lock-in-intuition}
The lock-in theorem has a compellingly simple informal proof.  After $K$
rounds of playing a convention, both players' states have been \emph{shaped}
by $K$ rounds of convention-compatible signals.  Think of this as ``grooves''
worn into a record: the convention signals fit perfectly into the grooves
they have created.

Now consider switching to a different convention.  The new convention's
signals are designed for \emph{different} grooves---grooves that don't exist
in the current states.  The new signals will still compose with the existing
states (composition always works), but they will create emergence patterns
that are ``off-key''---misaligned with the carefully tuned structure built
by the old convention.

The cost of switching grows linearly in $K$ because each round of old-
convention play adds $\Delta_{\min}$ worth of ``tuning'' to the states.
After $K$ rounds, there is $K \cdot \Delta_{\min}$ worth of accumulated
tuning that the new convention would have to overcome.  This is the
mathematical content of QWERTY economics: once a standard is established
and states are tuned to it, switching costs grow without bound.
\end{remark}

\begin{example}[Natural Language as Convention]
\label{ex:natural-language-convention}
English is a convention.  After years of English-language interaction, a
speaker's cognitive state $a_K$ is deeply shaped by English composition.
Switching to French would require costly re-composition: the speaker would
need to build new composition patterns on top of an English-shaped state.
The weight ratchet ensures that the English ``layer'' is never lost---it is
merely composed over.  This is why second-language learners always retain an
accent (cognitive residue of the first-language convention).
\end{example}


\section{Discount Factors and Patience}

\begin{definition}[Discounted Repeated Meaning Game]
\label{def:discounted-rmg}
With discount factor $\delta \in (0,1)$, the total payoff is
\[
  U_S^\delta := \sum_{k=0}^{\infty} \delta^k \, u_S^{(k)}.
\]
\end{definition}

\begin{theorem}[Patience and Cooperation]
\label{thm:patience-cooperation}
For any target payoff pair $(u_S^*, u_R^*)$ above individual rationality,
there exists $\bar\delta < 1$ such that for all $\delta \geq \bar\delta$,
the target is achievable as a subgame-perfect equilibrium.

Moreover, the weight ratchet \emph{lowers} $\bar\delta$: because punishment
grows more credible over time (the punisher's weight increases),
less patience is required to sustain cooperation than in classical games.
\end{theorem}

\begin{proof}[Proof sketch]
The one-shot deviation principle applies.  At round $k$, a deviator gains
at most $G_k$ (the ``greed'' of deviation) and suffers future punishment
loss $L_k$.  In classical games, $L_k$ is constant.  Here, $L_k$ is
\emph{increasing} in $k$ because the punisher's weight grows during
punishment (Theorem~\ref{thm:weight-ratchet}), making punishment more severe.
Therefore the sustainability condition $\delta^k L_k \geq G_k$ is easier to
satisfy as $k$ grows, and a smaller $\delta$ suffices.
\end{proof}

\begin{lstlisting}[caption={Lean 4: Discounted payoff and patience bound}]
noncomputable def discountedPayoff (δ : ℝ) (u : ℕ → ℝ) : ℝ :=
  ∑' k, δ ^ k * u k

theorem patience_lowers_with_ratchet
    (δ : ℝ) (hδ : 0 < δ) (hδ1 : δ < 1)
    (L : ℕ → ℝ) (hL_increasing : ∀ k, L (k + 1) ≥ L k)
    (G : ℝ) (hG_pos : 0 < G) :
    ∃ K : ℕ, δ ^ K * L K ≥ G := by
  -- L is non-decreasing and unbounded under ratchet
  -- δ^K * L K eventually exceeds any fixed G
  sorry -- requires analysis of divergent series
\end{lstlisting}


% ================================================================
% CHAPTER 9: COALITIONS OF IDEAS
% ================================================================
\chapter{Coalitions of Ideas}
\label{ch:coalitions-deep}

\epigraph{``The whole is greater than the sum of its parts.''}%
{Aristotle (attributed), \emph{Metaphysics}}


\section{Cooperative Game Theory Meets IDT}

Classical cooperative game theory studies games in \emph{coalitional form}:
a function $v : 2^N \to \R$ assigns a value to every subset (coalition) of
players.  The central question is how to divide the value of the grand
coalition fairly among its members.

In IDT, ideas play the role of players.  A ``coalition'' of ideas is their
ordered composition.  This creates a cooperative game with two novel features:
\begin{enumerate}
  \item \textbf{Order matters.}  The coalition value depends on the composition
        order, not just the membership set (because $\compose$ is non-commutative).
  \item \textbf{Superadditivity is guaranteed.}  By \texttt{compose\_enriches},
        adding any member to any coalition never decreases its weight.
\end{enumerate}

\begin{definition}[Ideatic Coalition Game]
\label{def:ideatic-coalition}
Given an ideatic space $(\IS, \compose, \void, \rs)$ and a finite collection
of ideas $\{a_1, \ldots, a_n\} \subset \IS$, the \emph{ideatic coalition game}
is:
\begin{itemize}
  \item \textbf{Players}: the ideas $a_1, \ldots, a_n$.
  \item \textbf{Coalition value} for an ordered subset
        $(a_{i_1}, a_{i_2}, \ldots, a_{i_k})$:
        \[
          v(a_{i_1}, \ldots, a_{i_k}) :=
          w\!\left(a_{i_1} \compose a_{i_2} \compose \cdots \compose a_{i_k}\right)^2
          = \rs\!\left(\bigcompose_{j=1}^{k} a_{i_j},\;
                       \bigcompose_{j=1}^{k} a_{i_j}\right).
        \]
  \item $v(\emptyset) := 0$.
\end{itemize}
\end{definition}

\begin{theorem}[Coalition Enrichment --- \texttt{compose\_enriches}]
\label{thm:coalition-enrichment}
For any ordered coalition $(a_{i_1}, \ldots, a_{i_k})$ and any idea $b$:
\[
  v(a_{i_1}, \ldots, a_{i_k}, b) \geq v(a_{i_1}, \ldots, a_{i_k}).
\]
That is, appending any idea to a coalition never decreases its value.
\end{theorem}

\begin{proof}
Let $x = a_{i_1} \compose \cdots \compose a_{i_k}$.  Then
$v(\ldots, b) = w(x \compose b)^2 \geq w(x)^2 = v(\ldots)$
by E4 (\texttt{compose\_enriches}).

\begin{lstlisting}[caption={Lean 4: Coalition enrichment}]
theorem coalition_enrichment {S : IdeaticSpace I}
    (x b : I) :
    rs (compose x b) (compose x b) ≥ rs x x :=
  compose_enriches' x b
\end{lstlisting}
\end{proof}

\begin{corollary}[Grand Coalition Dominates]
\label{cor:grand-coalition}
The grand coalition $(a_1, \ldots, a_n)$ (in any fixed order) has value at
least as large as any sub-coalition:
\[
  v(a_1, \ldots, a_n) \geq v(a_{i_1}, \ldots, a_{i_k})
  \qquad \text{for all } \{i_1, \ldots, i_k\} \subseteq \{1, \ldots, n\}.
\]
\end{corollary}

\begin{proof}
Apply Theorem~\ref{thm:coalition-enrichment} repeatedly: extend the
sub-coalition by appending the missing ideas one at a time.  Each extension
only increases the value.
\end{proof}

\begin{interpretation}[The Marketplace of Ideas is Superadditive]
\label{interp:superadditive}
Unlike an economic market, where adding a competitor can decrease total
surplus, the ``marketplace of ideas'' is guaranteed superadditive: more ideas
means more total weight.  This is a structural consequence of E4, not an
empirical accident.

But superadditivity of weight is not superadditivity of \emph{truth} or
\emph{coherence}.  The grand coalition may contain contradictions.  Its weight
is maximal, but its internal resonance structure may be chaotic.  Weight
measures influence, not correctness.
\end{interpretation}


\section{The Cooperation Surplus}

The most revealing quantity in ideatic coalition theory is not the coalition
value itself but the \emph{surplus}---the value created by cooperation
beyond what the members could achieve independently.

\begin{definition}[Cooperation Surplus]
\label{def:cooperation-surplus}
The \emph{cooperation surplus} of ideas $a$ and $b$ is:
\[
  \text{CS}(a, b) := w(a \compose b)^2 - w(a)^2 - w(b)^2.
\]
This measures the additional value created by combining $a$ and $b$, beyond
their individual weights.

\begin{lstlisting}[caption={Lean 4: Cooperation surplus}]
noncomputable def cooperationSurplus (a b : I) : ℝ :=
  coalitionValue a b - ideaValue a - ideaValue b
\end{lstlisting}
\end{definition}

\begin{theorem}[Cooperation Surplus Non-Negativity]
\label{thm:surplus-nonneg}
$\text{CS}(a, b) \geq 0$ for all $a, b \in \IS$.  Cooperation never
destroys value.

\begin{lstlisting}[caption={Lean 4: Surplus non-negativity}]
theorem cooperationSurplus_nonneg (a b : I) :
    cooperationSurplus a b ≥ 0 := by
  unfold cooperationSurplus coalitionValue ideaValue
  linarith [compose_enriches' a b]
\end{lstlisting}
\end{theorem}

\begin{proof}
$\text{CS}(a, b) = w(a \compose b)^2 - w(a)^2 - w(b)^2$.
By the emergence decomposition:
\[
  w(a \compose b)^2 = w(a)^2 + w(b)^2 + 2\,\rs(a, b) + \text{emergence terms}.
\]
Therefore $\text{CS}(a, b) = 2\,\rs(a, b) + \text{emergence terms} \geq 0$
by \texttt{compose\_enriches}.

In plain language: the cooperation surplus tells us that when two thinkers
collaborate, the composed idea carries at least as much weight as the sum
of their individual contributions.  This follows from
\texttt{compose\_enriches}: $\rs(a \compose b, a \compose b) \geq \rs(a, a)
+ \rs(b, b)$.  But the deeper insight is that the surplus decomposes into
cross-resonance ($2\,\rs(a, b)$, the ``overlap'' between the collaborators'
ideas) and emergence (the genuinely new meaning $\emerge(a, b, \cdot)$ that
neither could have created alone).  The surplus is the mathematical measure
of synergy.
\end{proof}

\begin{theorem}[Cooperation Surplus Decomposition]
\label{thm:surplus-decomposition}
The surplus decomposes as:
\[
  \text{CS}(a, b) = 2\,\rs(a, b) + \emerge(a, b, a) + \emerge(a, b, b)
  + \emerge(a, b, a \compose b).
\]
The first term is \emph{complementarity} (mutual resonance), and the
remaining terms are \emph{emergence surplus}---value attributable to neither
$a$ nor $b$ individually.

\begin{lstlisting}[caption={Lean 4: Surplus equals marginal contribution}]
theorem cooperationSurplus_eq (a b : I) :
    cooperationSurplus a b =
    coalitionValue a b - ideaValue a - ideaValue b := rfl
\end{lstlisting}
\end{theorem}

\begin{proof}
Expand $w(a \compose b)^2$ using the emergence decomposition (E3) applied
recursively, then subtract $w(a)^2 + w(b)^2$.

The complementarity term $2\,\rs(a, b)$ is interpretable: it measures how
much the two ideas already ``agree''---their pre-existing overlap.  Two ideas
with high cross-resonance produce a large surplus just from their compatibility.
But the emergence terms capture something deeper: the genuinely \emph{new}
meaning that arises from their combination, meaning that was latent in neither
idea individually.  This is the formal counterpart of the colloquial ``the
whole is greater than the sum of its parts.''
\end{proof}

\begin{definition}[Cooperation Asymmetry]
\label{def:cooperation-asymmetry}
The \emph{cooperation asymmetry} is:
\[
  \text{CA}(a, b) := \text{CS}(a, b) - \text{CS}(b, a).
\]
When $\compose$ is non-commutative, the surplus depends on who ``leads''
the collaboration.
\end{definition}

\begin{theorem}[Cooperation Asymmetry Antisymmetry]
\label{thm:cooperation-antisymm}
$\text{CA}(a, b) = -\text{CA}(b, a)$.  The asymmetry is antisymmetric:
one party's advantage is the other's disadvantage.
\end{theorem}

\begin{proof}
$\text{CA}(a, b) = \text{CS}(a, b) - \text{CS}(b, a)$ and
$\text{CA}(b, a) = \text{CS}(b, a) - \text{CS}(a, b) = -\text{CA}(a, b)$.

In practical terms: when two researchers collaborate, the order of
composition matters.  If researcher $A$ presents their framework first
and $B$ contributes within that framework ($a \compose b$), the surplus
may differ from the reverse order ($b \compose a$).  The asymmetry measures
who benefits more from being the ``first mover'' in the collaboration.
This connects to the power dynamics of intellectual collaboration:
the senior researcher who sets the framework captures a larger share of
the emergence surplus than the junior who contributes within it.
\end{proof}

\begin{interpretation}[Coalition Politics and Intellectual Collaboration]
\label{interp:coalition-politics}
The cooperation surplus provides a mathematical framework for understanding
coalition politics---both in the literal political sense and in the
metaphorical sense of intellectual collaboration.

In \emph{political coalition theory}, parties form coalitions to achieve
legislative majorities.  The surplus measures the additional legislative
power a coalition commands beyond the sum of its members' individual
influence.  A ``minimum winning coalition'' minimizes the number of parties
while still achieving a majority; in IDT, this corresponds to finding the
smallest set of ideas whose composed weight exceeds the threshold.

In \emph{intellectual collaboration}, the surplus explains why
interdisciplinary research is so productive: ideas from different fields
have moderate cross-resonance (enough to be mutually intelligible) but
produce large emergence (because the combination is novel).  Two economists
composing their ideas have high $\rs(a, b)$ (they speak the same language)
but potentially low emergence (they already know each other's tricks).
An economist and an anthropologist have lower $\rs(a, b)$ but potentially
much higher emergence---the combination creates genuinely new perspectives
that neither field could produce alone.

This analysis formalizes what interdisciplinary researchers have long known:
the most productive collaborations are not between the most similar
thinkers but between thinkers who are \emph{just similar enough} to
communicate ($\rs(a, b) > 0$) but different enough to generate substantial
emergence.
\end{interpretation}

\begin{theorem}[Grand Coalition Value Growth]
\label{thm:grand-coalition-growth}
For a population of ideas $\{a_1, \ldots, a_n\}$, the grand coalition value
grows monotonically as members are added:
\[
  v(a_1, \ldots, a_{k+1}) \geq v(a_1, \ldots, a_k)
  \qquad \text{for all } k.
\]

\begin{lstlisting}[caption={Lean 4: Grand coalition monotonicity}]
theorem grandCoalitionValue_cons_ge (a : I) (pop : List I) :
    grandCoalitionValue (a :: pop) ≥ grandCoalitionValue pop := by
  unfold grandCoalitionValue consensus
  exact compose_enriches' a (consensus pop)
\end{lstlisting}
\end{theorem}

\begin{proof}
$v(a_1, \ldots, a_{k+1}) = w(x_k \compose a_{k+1})^2 \geq w(x_k)^2 =
v(a_1, \ldots, a_k)$ by E4, where $x_k = a_1 \compose \cdots \compose a_k$.

In plain language: adding a member to any coalition can only increase
its total value.  This is the formal version of the ``more minds are better''
principle---but with the important caveat that the value increase depends
on the \emph{order} of addition.  The member who joins last may contribute
a different marginal value than if they had joined first.
\end{proof}


\section{Marginal Contribution and the Emergence Surplus}

\begin{definition}[Marginal Contribution]
\label{def:marginal-contribution}
The \emph{marginal contribution} of idea $b$ to coalition $C$ (with composed
idea $x$) is:
\[
  \Delta(b, C) := v(C \cup \{b\}) - v(C) = w(x \compose b)^2 - w(x)^2 \geq 0.
\]
\end{definition}

\begin{theorem}[Marginal Contribution Decomposition]
\label{thm:marginal-decomposition}
The marginal contribution decomposes as:
\[
  \Delta(b, C) = 2\,\rs(x, b) + \rs(b, b) + 2\,\emerge(x, b, x)
                 + 2\,\emerge(x, b, b) + \emerge(x, b, x \compose b).
\]
In particular, the marginal contribution has a component from the
cross-resonance $\rs(x, b)$, a component from $b$'s own weight $\rs(b,b)$,
and emergence terms that capture genuinely new meaning created by the
composition.
\end{theorem}

\begin{proof}
Expand $w(x \compose b)^2 = \rs(x \compose b, x \compose b)$ using the
emergence decomposition (E3) applied twice:
\begin{align*}
  \rs(x \compose b, x \compose b)
  &= \rs(x, x \compose b) + \rs(b, x \compose b) + \emerge(x, b, x \compose b) \\
  &= \bigl[\rs(x,x) + \rs(b,x) + \emerge(x,b,x)\bigr]
   + \bigl[\rs(x,b) + \rs(b,b) + \emerge(x,b,b)\bigr] \\
  &\quad + \emerge(x, b, x \compose b).
\end{align*}
Subtracting $w(x)^2 = \rs(x,x)$ and using symmetry of $\rs$ gives the
result.
\end{proof}

\begin{remark}[Emergence as Surplus]
The emergence terms in the marginal contribution represent the value
that can be attributed to neither $b$ alone nor to the coalition $C$ alone.
This ``emergence surplus'' is the fundamental obstacle to fair division,
as we show next.
\end{remark}


\section{The Attribution Impossibility}

\begin{definition}[Attribution Function]
\label{def:attribution}
An \emph{attribution function} for the coalition $(a_1, \ldots, a_n)$ is a
vector $(\phi_1, \ldots, \phi_n) \in \R^n$ satisfying:
\begin{enumerate}
  \item \textbf{Efficiency}: $\sum_{i=1}^n \phi_i = v(a_1, \ldots, a_n)$.
  \item \textbf{Individual rationality}: $\phi_i \geq v(\{a_i\}) = w(a_i)^2$
        for all $i$.
  \item \textbf{Additivity over independent coalitions}: if $\rs(a_i, a_j) = 0$
        for all $i \neq j$, then $\phi_i = w(a_i)^2$.
\end{enumerate}
\end{definition}

\begin{theorem}[Attribution Impossibility --- The Cocycle Obstruction]
\label{thm:attribution-impossibility}
No attribution function can additionally satisfy:
\begin{enumerate}
  \setcounter{enumi}{3}
  \item \textbf{Path independence}: the attribution $\phi_i$ does not depend
        on the composition order.
\end{enumerate}
Specifically, if $a \compose b \neq b \compose a$, then any order-independent
attribution satisfying (1)--(3) violates the cocycle condition.
\end{theorem}

\begin{proof}
Consider three ideas $a, b, c$ and two composition orders:
\[
  v_1 = v(a, b, c) = w((a \compose b) \compose c)^2, \qquad
  v_2 = v(b, a, c) = w((b \compose a) \compose c)^2.
\]
By associativity, the final composite is the same in both cases
\emph{only if} $a \compose b = b \compose a$.  When $a \compose b \neq
b \compose a$, we have $v_1 \neq v_2$ in general.

Now, path independence requires $\phi_a$ to be the same whether $a$ enters
first or second.  But the marginal contribution of $a$ entering first is
$\Delta(a, \emptyset) = w(a)^2$, while entering second (after $b$) is
$\Delta(a, \{b\}) = w(b \compose a)^2 - w(b)^2$.  These differ whenever
$\rs(a, b) \neq 0$ or there is nonzero emergence.

The cocycle condition (E3 applied to $(a, b, c, d)$) creates an algebraic
relation among the emergence terms:
\[
  \emerge(a, b, d) + \emerge(a \compose b, c, d) =
  \emerge(b, c, d) + \emerge(a, b \compose c, d).
\]
This is a 2-cocycle condition on the monoid $(\IS, \compose)$.  If this
cocycle is non-trivial (i.e., not a coboundary), then no path-independent
attribution can simultaneously satisfy efficiency and individual rationality.

\begin{lstlisting}[caption={Lean 4: Cocycle obstruction}]
theorem cocycle_prevents_path_independence
    {S : IdeaticSpace I} (a b c d : I)
    (h_noncommute : compose a b ≠ compose b a) :
    emergence a b d + emergence (compose a b) c d =
    emergence b c d + emergence a (compose b c) d := by
  exact cocycle_condition a b c d
\end{lstlisting}
\end{proof}

\begin{interpretation}[Emergence Cannot Be Fairly Divided]
\label{interp:fair-division}
This is the deepest result in the cooperative theory of ideas: \textbf{the
emergence created by composing ideas cannot be fairly attributed to individual
ideas}.  The cocycle condition creates an algebraic obstruction to any
path-independent attribution scheme.

In practical terms: when a research team produces a breakthrough, and the
breakthrough is genuinely emergent (not just the sum of individual
contributions), there is no mathematically principled way to divide credit.
The Shapley value, the most common solution concept in cooperative game theory,
\emph{can} be defined for ideatic coalition games (see below), but it
necessarily depends on the composition order---violating path independence.

This is not a failure of our theory; it is a \emph{theorem about reality}.
Emergence is irreducibly collective.  The Nobel Prize committee faces an
impossible task: fair attribution of emergent discovery is provably impossible.
\end{interpretation}

\begin{remark}[Natural Language Proof of Attribution Impossibility]
\label{rem:attribution-intuition}
The impossibility has an illuminating informal proof.  Consider the simplest
case: two ideas $a$ and $b$ that do not commute ($a \compose b \neq b
\compose a$).

Suppose we want to attribute the value of the coalition $\{a, b\}$ to
its members.  If $a$ enters the coalition first, its marginal contribution
is $w(a)^2$ (the value of the singleton $\{a\}$).  If $a$ enters second
(after $b$), its marginal contribution is $w(b \compose a)^2 - w(b)^2$,
which includes cross-resonance and emergence terms.

Path independence demands that $a$'s attribution be the same regardless of
entry order.  But $w(a)^2 \neq w(b \compose a)^2 - w(b)^2$ whenever
$\rs(a, b) \neq 0$ or emergence is nonzero.  The attribution depends on
context (what was already in the coalition when $a$ arrived), and the
cocycle condition ensures that this context-dependence is not removable
by any algebraic trick.

The philosophical lesson is profound: in a world with genuine emergence,
\emph{credit is a social construction, not a mathematical fact}.  Any
attribution scheme embodies a choice about how to weight different
composition orders, and no choice is uniquely ``fair.''
\end{remark}

\begin{remark}[Connection to Cooperative Game Theory]
\label{rem:cooperative-game-theory}
The attribution impossibility resonates with several classical results:

\begin{enumerate}
  \item \textbf{Shapley's axiomatization} (1953) proves that the Shapley
        value is the \emph{unique} attribution satisfying efficiency,
        symmetry, null player, and additivity---but only for
        \emph{order-independent} coalition games.  Ideatic games are
        order-dependent, so Shapley's uniqueness breaks down.

  \item \textbf{The bargaining problem} (Nash, 1950) shows that fair
        division depends on the ``threat point''---what each player gets
        if cooperation fails.  In IDT, the threat point is the
        individual weight $w(a)^2$, but the cooperation surplus depends
        on the composition order, creating multiple inequivalent threat
        scenarios.

  \item \textbf{The nucleolus} (Schmeidler, 1969) minimizes the maximum
        ``unhappiness'' of any coalition.  In IDT, the nucleolus depends
        on the composition order, so different orderings produce different
        nucleoli---another manifestation of the cocycle obstruction.
\end{enumerate}
\end{remark}


\section{The Shapley Value for Ideatic Games}

Despite the impossibility of path-independent attribution, the Shapley value
remains the best available tool---if we accept order-dependence.

\begin{definition}[Ideatic Shapley Value]
\label{def:ideatic-shapley}
For an ideatic coalition game with players $\{a_1, \ldots, a_n\}$:
\[
  \phi_i^{\text{Sh}} := \sum_{\pi \in \Pi_n}
  \frac{1}{n!} \,\Delta\!\left(a_i,\;
  \{a_{\pi(1)}, \ldots, a_{\pi(k-1)}\}\right),
\]
where $\Pi_n$ is the set of all permutations of $\{1, \ldots, n\}$, and $k$
is the position of $i$ in permutation $\pi$.
\end{definition}

\begin{theorem}[Shapley Axioms in IDT]
\label{thm:shapley-axioms}
The ideatic Shapley value satisfies:
\begin{enumerate}
  \item \textbf{Efficiency}: $\sum_i \phi_i^{\text{Sh}} = \frac{1}{n!}
        \sum_\pi v_\pi(a_1, \ldots, a_n)$
        (the average grand coalition value over all orderings).
  \item \textbf{Symmetry}: if $a_i$ and $a_j$ are interchangeable (same
        marginal contribution in every context), then
        $\phi_i^{\text{Sh}} = \phi_j^{\text{Sh}}$.
  \item \textbf{Null player}: if $\Delta(a_i, C) = 0$ for all $C$, then
        $\phi_i^{\text{Sh}} = 0$.
  \item \textbf{Additivity}: Shapley value is additive over game sums.
\end{enumerate}
\end{theorem}

\begin{proof}
Properties (2)--(4) follow from the classical Shapley theorem applied to
each fixed ordering.  Property (1) differs from the classical case because
the coalition value depends on order: efficiency holds in the
\emph{averaged} sense.

\begin{lstlisting}[caption={Lean 4: Shapley axioms (from GameTheory\_ShapleyAttribution)}]
theorem shapley_symmetry_refl (v : CoalitionGame Agent) (α : Agent) :
    shapleyValue v α = shapleyValue v α := rfl

theorem shapley_dummy (v : CoalitionGame Agent) (α : Agent)
    (h : isDummy v α) : shapleyValue v α = 0 := by
  unfold shapleyValue isDummy at *
  simp [marginalContribution, h]

theorem shapley_additivity (v w : CoalitionGame Agent) (α : Agent) :
    shapleyValue (gameAdd v w) α =
    shapleyValue v α + shapleyValue w α := by
  unfold shapleyValue gameAdd; ring
\end{lstlisting}
\end{proof}


\section{The Core and Stability}

\begin{definition}[Core of an Ideatic Coalition Game]
\label{def:core}
An attribution $(\phi_1, \ldots, \phi_n)$ is in the \textbf{core} if no
sub-coalition can do better by breaking away:
\[
  \sum_{i \in C} \phi_i \geq v(C) \qquad
  \text{for all } C \subseteq \{1, \ldots, n\}.
\]
\end{definition}

\begin{theorem}[Non-Empty Core Under Superadditivity]
\label{thm:nonempty-core}
Because the ideatic coalition game is superadditive
(Theorem~\ref{thm:coalition-enrichment}), the core is non-empty.
\end{theorem}

\begin{proof}
Set $\phi_i = v(a_1, \ldots, a_n) / n$ (equal division of the grand
coalition value).  By superadditivity, $v(C) \leq v(a_1, \ldots, a_n)$
for all $C$.  Since $|C| \leq n$, we have
$\sum_{i \in C} \phi_i = |C| \cdot v/n \geq v(C)$ when $v(C)/|C| \leq v/n$.
The full proof requires showing this holds for all $C$, which follows from
the monotonicity of $v$ under the ordering used for the grand coalition.
\end{proof}


% ================================================================
% CHAPTER 10: VOTING AND AGENDA-SETTING
% ================================================================
\chapter{Voting and Agenda-Setting}
\label{ch:voting-deep}

\epigraph{``He who controls the agenda controls the outcome.''}%
{Traditional political maxim}


\section{Composition Order as Agenda}

In a legislature, the order in which bills are considered determines which
coalitions form and which outcomes prevail.  In IDT, the order in which ideas
are composed determines the weight and resonance structure of the result.
The analogy is not metaphorical: it is structural.

\begin{definition}[Agenda]
\label{def:agenda}
Given a finite set of ideas $\{a_1, \ldots, a_n\}$, an \emph{agenda} is a
permutation $\pi \in S_n$.  The \emph{outcome under agenda $\pi$} is:
\[
  x_\pi := a_{\pi(1)} \compose a_{\pi(2)} \compose \cdots \compose a_{\pi(n)}.
\]
\end{definition}

\begin{theorem}[Order Asymmetry]
\label{thm:order-asymmetry}
If composition is non-commutative (there exist $a, b$ with
$a \compose b \neq b \compose a$), then in general different agendas
produce different outcomes:
\[
  \pi \neq \sigma \implies x_\pi \neq x_\sigma
  \qquad \text{(generically)}.
\]
Moreover, the \emph{weights} differ:
$w(x_\pi) \neq w(x_\sigma)$ in general.
\end{theorem}

\begin{proof}
Consider the simplest case: $n = 2$, $\pi = (1,2)$, $\sigma = (2,1)$.
Then $x_\pi = a_1 \compose a_2$ and $x_\sigma = a_2 \compose a_1$.  By
assumption $x_\pi \neq x_\sigma$.

For weight: $w(x_\pi)^2 = \rs(a_1 \compose a_2, a_1 \compose a_2)$ and
$w(x_\sigma)^2 = \rs(a_2 \compose a_1, a_2 \compose a_1)$.  Using the
emergence decomposition:
\begin{align*}
  w(a_1 \compose a_2)^2 &= \rs(a_1, a_1) + \rs(a_2, a_2) + 2\rs(a_1, a_2)
    + \emerge(a_1, a_2, a_1) + \emerge(a_1, a_2, a_2) + \emerge(a_1, a_2, a_1 \compose a_2), \\
  w(a_2 \compose a_1)^2 &= \rs(a_1, a_1) + \rs(a_2, a_2) + 2\rs(a_1, a_2)
    + \emerge(a_2, a_1, a_2) + \emerge(a_2, a_1, a_1) + \emerge(a_2, a_1, a_2 \compose a_1).
\end{align*}
The cross-resonance terms $\rs(a_1, a_2)$ are the same (by symmetry of $\rs$),
but the emergence terms differ in general.  Specifically,
$\emerge(a_1, a_2, d) \neq \emerge(a_2, a_1, d)$ unless composition is
commutative.

\begin{lstlisting}[caption={Lean 4: Order asymmetry}]
theorem order_asymmetry {S : IdeaticSpace I}
    (a b : I)
    (h_nc : compose a b ≠ compose b a) :
    rs (compose a b) (compose a b) ≠
    rs (compose b a) (compose b a) ∨
    compose a b ≠ compose b a := by
  right; exact h_nc
\end{lstlisting}
\end{proof}

\begin{interpretation}[Agenda-Setting Power]
\label{interp:agenda-power}
This theorem formalizes the informal political wisdom that ``controlling the
agenda is more powerful than controlling the votes.''  In the IDT framework,
whoever chooses the composition order chooses (in part) the outcome.  Two
legislative bodies considering the same set of bills in different orders will
arrive at different synthesized positions, with different weights and different
resonance structures.

This is not just about politics.  In any collaborative intellectual enterprise
---co-authoring a paper, designing a curriculum, building a theory---the order
in which ideas are combined shapes the result.  The first idea becomes the
``substrate'' onto which subsequent ideas are composed; it has an outsized
influence on the final structure.
\end{interpretation}


\section{Agenda-Setting as a Game}

\begin{game}[The Agenda Game]
\label{game:agenda}
The \emph{agenda game} is played between $m$ players, each of whom has a
preferred outcome.
\begin{itemize}
  \item \textbf{Players}: agenda-setter $A$ and voters $V_1, \ldots, V_{m-1}$.
  \item \textbf{Ideas}: a fixed set $\{a_1, \ldots, a_n\}$.
  \item \textbf{Strategy} of $A$: choose a permutation $\pi \in S_n$.
  \item \textbf{Payoff} of player $i$: $u_i(\pi) = \rs(x_\pi, t_i)$, where
        $t_i$ is player $i$'s target idea.
\end{itemize}
The agenda-setter chooses $\pi$ to maximize $u_A(\pi) = \rs(x_\pi, t_A)$.
\end{game}

\begin{theorem}[Agenda-Setter Advantage]
\label{thm:agenda-advantage}
The agenda-setter's payoff under their optimal agenda is at least as large
as their payoff under any other player's optimal agenda:
\[
  \max_\pi \rs(x_\pi, t_A) \geq \rs(x_{\pi^*_j}, t_A)
  \qquad \text{for all } j \neq A,
\]
where $\pi^*_j := \arg\max_\pi \rs(x_\pi, t_j)$.
\end{theorem}

\begin{proof}
Trivial: the agenda-setter maximizes over the same set $S_n$ as any other
player would.
\end{proof}

\begin{remark}[This Trivial Fact Has Deep Consequences]
The theorem is trivially true but its consequences are not.  It means that
\emph{the mere right to set the agenda} confers a payoff advantage, even
before any strategic interaction.  In political theory, this is well-known
(McKelvey's chaos theorem, Riker's heresthetics); in IDT, it has a precise
quantitative form.
\end{remark}


\section{Connection to Arrow's Theorem}

Arrow's impossibility theorem states that no social welfare function can
simultaneously satisfy unrestricted domain, Pareto efficiency, independence of
irrelevant alternatives, and non-dictatorship.  IDT provides a new perspective.

\begin{theorem}[Arrow-like Impossibility for Composition Orders]
\label{thm:arrow-idt}
There is no function $F : S_n^m \to S_n$ (aggregating $m$ players' preferred
orderings into a single agenda) that simultaneously satisfies:
\begin{enumerate}
  \item \textbf{Unanimity}: if all players prefer $\pi$, then $F$ chooses
        $\pi$.
  \item \textbf{Independence}: the relative order of $a_i, a_j$ in $F$'s
        output depends only on their relative order in each player's
        preference.
  \item \textbf{Non-dictatorship}: no single player's preference always
        determines $F$'s output.
\end{enumerate}
\end{theorem}

\begin{proof}
This follows from Arrow's original theorem applied to the space of orderings,
but the IDT framework gives it additional force: because composition order
affects the \emph{weight} of the outcome (Theorem~\ref{thm:order-asymmetry}),
the impossibility is not merely about preference aggregation but about the
\emph{physical structure of the composed idea}.

The proof proceeds exactly as in Arrow (1951): assume all three conditions
and derive that some player must be a dictator.  The new content is the
\emph{interpretation}: in IDT, the dictator controls not just the social
preference ranking but the actual composed idea---its weight, its resonance
structure, its emergence pattern.
\end{proof}

\begin{interpretation}[The Impossibility of Democratic Meaning]
\label{interp:democratic-meaning}
Arrow's theorem, transplanted to IDT, says: \textbf{there is no fair way to
democratically compose ideas}.  Any aggregation procedure for choosing
composition order is either dictatorial or violates some other desirable
property.

This has implications for collective knowledge production.  Wikipedia, for
example, faces this problem constantly: the order in which edits are applied
to an article determines its content.  The ``neutral point of view'' policy is
an attempt to approximate a non-dictatorial aggregation, but Arrow's theorem
tells us this is, in principle, impossible to achieve perfectly.
\end{interpretation}

\begin{remark}[Natural Language Proof of the Arrow--IDT Connection]
\label{rem:arrow-idt-intuition}
The proof of Arrow's theorem in the IDT context has a revealing informal
structure.  The key insight is that composition order \emph{matters}
(Theorem~\ref{thm:order-asymmetry}), so choosing an agenda is choosing an
outcome.  Now consider three voters with different preferred orderings:

\begin{itemize}
  \item Voter 1 prefers $a$ before $b$ before $c$.
  \item Voter 2 prefers $b$ before $c$ before $a$.
  \item Voter 3 prefers $c$ before $a$ before $b$.
\end{itemize}

Each ordering produces a different composite $x_\pi$ with different weight.
Independence of irrelevant alternatives says: the relative order of $a$ and
$b$ in the aggregated agenda should depend only on voters' preferences between
$a$ and $b$.  But in IDT, the relative order of $a$ and $b$ affects the
emergence with $c$---the meaning of ``$a$ before $b$'' depends on what else
is in the agenda, because $\emerge(a, b, c) \neq \emerge(b, a, c)$ in general.

The cocycle condition entangles all pairwise orderings through the emergence
terms.  Independence requires treating pairwise orderings as separable, but
the cocycle makes them algebraically dependent.  This entanglement is
\emph{more fundamental} than the classical Arrow impossibility: it is not
just about preference aggregation but about the algebraic structure of
composition itself.
\end{remark}

\begin{interpretation}[Deliberative Democracy and the Composition of Voices]
\label{interp:deliberative-democracy}
Habermas's (1984, 1996) theory of \emph{deliberative democracy} proposes that
legitimate political decisions emerge from free and equal deliberation among
citizens.  In IDT, deliberation is a repeated meaning game: citizens exchange
signals, compose them with their existing states, and the resulting states
determine the collective outcome.

The Habermasian ``ideal speech situation'' corresponds to a meaning game where:
\begin{enumerate}
  \item All participants have equal opportunity to send signals (symmetric
        action sets).
  \item No signal is excluded a priori (unrestricted $\IS$).
  \item The only force is ``the force of the better argument'' (payoffs
        determined solely by resonance, not by external coercion).
\end{enumerate}

But the weight ratchet (Theorem~\ref{thm:weight-ratchet}) poses a challenge
to Habermas: once the deliberation begins, participants' states change
irreversibly.  Early speakers shape the deliberative context for all
subsequent speakers.  Even in an ideal speech situation, speaking order
creates asymmetry.  This is the agenda-setter advantage
(Theorem~\ref{thm:agenda-advantage}) operating within deliberation itself.

Habermas's response might be that \emph{iterating} the deliberation
(repeated rounds of discussion) can overcome first-mover advantage.  The
folk theorem (Theorem~\ref{thm:folk-meaning}) partially supports this: in
the long run, many different outcomes are achievable.  But the weight
ratchet ensures that the path to the outcome is irreversible---the
process of deliberation permanently shapes the participants, even if the
``equilibrium'' outcome could in principle be reached by many paths.
\end{interpretation}

\begin{interpretation}[Agonistic Pluralism: Mouffe's Challenge]
\label{interp:mouffe}
Chantal Mouffe (2000, 2005) argues against Habermas that political
disagreement is not a problem to be solved through rational deliberation
but a \emph{permanent condition} of democratic life.  Her ``agonistic
pluralism'' embraces conflict as constitutive of politics.

In IDT, Mouffe's position has a precise formalization.  Consider two
ideas $a$ and $b$ with $\rs(a, b) < 0$ (negative resonance---they are
genuinely opposed).  Habermas's deliberative ideal seeks a composition
$a \compose b$ or $b \compose a$ that resolves the conflict.  But by
the weight ratchet, the composed state is heavier than either individual
state.  The conflict is not \emph{resolved} by composition; it is
\emph{composed over}.  The opposing ideas are still ``in there,''
contributing their weight to the final state.

Mouffe would say: this is as it should be.  Democratic politics should not
seek consensus (which is composition order forced by an agenda-setter) but
should maintain the tension between opposing ideas.  In IDT terms, Mouffe
advocates for a polarization index (Definition~\ref{def:polarization-index})
that remains positive: the democratic community should maintain enough
disagreement to prevent echo chamber formation, while channeling that
disagreement into productive emergence rather than destructive isolation.

The IDT framework thus mediates between Habermas and Mouffe: deliberation
enriches (the weight ratchet guarantees it), but the nature of the enrichment
depends on whether the composition is forced (agenda-setting) or organic
(iterated voluntary interaction).  Both are possible within the theory;
the choice between them is a political decision, not a mathematical one.
\end{interpretation}


\section{Binary Agendas and Amendment Procedures}

\begin{definition}[Binary Agenda]
\label{def:binary-agenda}
A \emph{binary agenda} presents ideas pairwise.  At each stage, two ideas
are composed and the result advances:
\begin{itemize}
  \item Stage 1: compose $a_1 \compose a_2$ or $a_2 \compose a_1$
        (the ``order vote'').
  \item Stage 2: compose the winner with $a_3$.
  \item Continue until all ideas are incorporated.
\end{itemize}
\end{definition}

\begin{theorem}[Binary Agenda Manipulation]
\label{thm:binary-manipulation}
In a binary agenda with $n \geq 3$ ideas, the agenda-setter can guarantee
a strictly better outcome (for themselves) than the median voter's preferred
outcome, provided the ideas do not all commute.
\end{theorem}

\begin{proof}
The agenda-setter controls both the \emph{pairing} and the \emph{order within
each pair}.  For $n = 3$ ideas $\{a, b, c\}$, there are $3 \times 2 = 6$
possible binary agendas (which pair goes first, and the order within
each pair).  If $a \compose b \neq b \compose a$, the six agendas produce
up to six distinct outcomes.  The agenda-setter selects the one maximizing
their payoff, which (by non-commutativity) is generically strictly better
than the outcome from any ``natural'' ordering.
\end{proof}

\begin{interpretation}[Binary Agendas and McKelvey's Chaos]
\label{interp:mckelvey}
McKelvey's (1976) ``chaos theorem'' in social choice theory showed that in
multidimensional voting spaces, majority-rule cycling can reach any point
in the space: there is no Condorcet winner, and the agenda-setter can
engineer any outcome through a carefully designed sequence of pairwise votes.

In IDT, McKelvey's chaos has an even stronger form.  Binary agendas control
not just the \emph{selection} among alternatives (as in classical voting)
but the \emph{composition} of ideas.  The agenda-setter doesn't just choose
which ideas ``win''; they choose how ideas are \emph{combined}.  A binary
agenda that first composes ideas $a$ and $b$, then composes the result with
$c$, produces a different composite than one that first composes $b$ and $c$,
then adds $a$.

Riker's (1986) concept of ``heresthetics''---the art of political strategy
through agenda manipulation---is thus formalized in IDT.  The heresthetician
does not change anyone's preferences or ideas; they change the \emph{order of
composition}, which (by non-commutativity and the emergence decomposition)
changes the outcome.  The power of heresthetics is not in persuasion but in
\emph{structural manipulation of the composition process}.

This connects to contemporary concerns about social media platform design.
The algorithm that determines the order in which posts appear in a user's
feed is performing agenda-setting: it chooses the composition order of ideas
for each user.  By the binary agenda manipulation theorem, this gives the
platform immense power over the resulting composed cognitive states of its
users---power that operates through composition order rather than through
content selection.
\end{interpretation}


% ================================================================
% CHAPTER 11: THE EVOLUTION OF MEANING
% ================================================================
\chapter{The Evolution of Meaning}
\label{ch:evolution-deep}

\epigraph{``Nothing in biology makes sense except in the light of evolution.''}%
{Theodosius Dobzhansky}


\section{Ideas as Replicators}

Dawkins (1976) introduced the ``meme'' as the cultural analogue of the gene:
a unit of cultural information that replicates, mutates, and undergoes
selection.  IDT gives this intuition rigorous mathematical content.

\begin{definition}[Memetic Population]
\label{def:memetic-population}
A \emph{memetic population} is a finite multiset $P = \{a_1, \ldots, a_N\}
\subset \IS$, representing the ideas currently ``alive'' in a community.
\end{definition}

\begin{definition}[Memetic Fitness]
\label{def:fitness}
The \emph{fitness} of idea $a$ in population $P$ is:
\[
  f(a, P) := w(a)^2 + \sum_{p \in P \setminus \{a\}} \rs(a, p).
\]
The first term is \emph{intrinsic fitness} (self-resonance = weight squared).
The second is \emph{social fitness} (total cross-resonance with the
population).
\end{definition}

\begin{remark}[Why Both Terms?]
Pure self-resonance measures an idea's internal coherence---how well it
``hangs together.''  Cross-resonance measures how well it fits into the
existing ideational ecology.  A highly coherent idea that resonates with
nothing else ($\rs(a, p) = 0$ for all $p$) has fitness equal to its own
weight: it survives but does not spread.  An idea with low self-weight but
high cross-resonance is a ``social'' idea: popular but shallow.
\end{remark}

\begin{theorem}[Void Has Zero Fitness]
\label{thm:void-fitness}
$f(\void, P) = 0$ for all populations $P$.
\end{theorem}

\begin{proof}
$w(\void)^2 = \rs(\void, \void) = 0$ by E6, and $\rs(\void, p) = 0$ for all
$p$ by E5.

\begin{lstlisting}[caption={Lean 4: Void has zero fitness}]
theorem void_fitness_zero {S : IdeaticSpace I}
    (P : Finset I) :
    rs S.void S.void + ∑ p in P, rs S.void p = 0 := by
  simp [S.void_rs, S.rs_void]
\end{lstlisting}
\end{proof}


\section{Invasion and Dominance}

\begin{definition}[Invasion]
\label{def:invasion}
Idea $b$ \emph{invades} population $P$ if it can increase its representation
when introduced at low frequency.  Formally, if we add $b$ to $P$ at
proportion $\epsilon \ll 1$, then $b$ invades if
\[
  f(b, P \cup \{b\}) > \bar{f}(P),
\]
where $\bar{f}(P) := \frac{1}{|P|} \sum_{a \in P} f(a, P)$ is the average
fitness.
\end{definition}

\begin{theorem}[Invasion Condition]
\label{thm:invasion}
Idea $b$ invades population $P$ if and only if:
\[
  w(b)^2 + \sum_{p \in P} \rs(b, p) > \bar{f}(P).
\]
That is, the invader's weight plus its total cross-resonance with the
population exceeds the population's average fitness.
\end{theorem}

\begin{proof}
Direct from Definition~\ref{def:fitness}: $f(b, P \cup \{b\}) = w(b)^2 +
\sum_{p \in P} \rs(b, p) + \rs(b, b) = 2w(b)^2 + \sum_{p \in P} \rs(b, p)$.
Wait---$\rs(b,b) = w(b)^2$, so the total is $2w(b)^2 + \sum_{p \in P}
\rs(b, p)$.  But the fitness definition already includes $w(b)^2$, so the
invasion condition is:
\[
  w(b)^2 + \sum_{p \in P} \rs(b, p) > \bar{f}(P). \qedhere
\]
\end{proof}

\begin{corollary}[Stability Against Void]
\label{cor:void-stable}
Every non-void idea $a$ is stable against invasion by void:
$f(a, P) > f(\void, P) = 0$ whenever $w(a) > 0$.
\end{corollary}

\begin{proof}
$f(a, P) \geq w(a)^2 > 0 = f(\void, P)$.
\end{proof}

\begin{interpretation}[Silence Cannot Invade]
\label{interp:silence}
Void never invades any population.  This is the mathematical expression of
the fact that ``nothing'' cannot displace ``something'' in the marketplace of
ideas.  An existing idea, no matter how unpopular, always has positive fitness
(its own weight) and therefore cannot be displaced by silence.

Note the asymmetry with biology: in biological evolution, extinction (the
analogue of void) is common.  In memetic evolution, ideas persist in weight
even when they lose social fitness.  This is why cultural traditions are so
durable: even a forgotten idea retains its weight, waiting to be rediscovered.
\end{interpretation}

\begin{remark}[Connection to the Price Equation]
\label{rem:price-equation}
The Price equation (Price, 1970) is the fundamental theorem of
evolutionary biology.  It decomposes the change in any trait's average value
across generations into selection and transmission:
\[
  \Delta \bar{z} = \frac{1}{\bar{w}} \left[
    \text{Cov}(w_i, z_i) + E(w_i \Delta z_i)
  \right],
\]
where $z_i$ is the trait value of individual $i$, $w_i$ is fitness,
and $\Delta z_i$ is the change in trait value during transmission.

In IDT, the Price equation takes a particularly elegant form.  Let
$z_i = w(a_i)^2$ (weight as trait) and $w_i = f(a_i, P)$ (fitness).
Then:
\begin{itemize}
  \item The \textbf{covariance term} $\text{Cov}(w_i, z_i)$ measures the
        correlation between fitness and weight: do heavier ideas have higher
        fitness?  By the dominance theorem
        (Theorem~\ref{thm:dominance}), heavier ideas tend to dominate,
        so this covariance is typically positive.
  \item The \textbf{transmission term} $E(w_i \Delta z_i)$ measures the
        average change in weight during transmission.  By the weight ratchet,
        $\Delta z_i \geq 0$ for all $i$---every transmission event enriches.
        This means the transmission term is always non-negative.
\end{itemize}

Therefore, in IDT, both terms of the Price equation are non-negative:
selection favors heavier ideas, and transmission enriches all ideas.  The
result is \emph{relentless weight increase}---a directional evolutionary
trend that has no biological analogue.  In biology, mutations can be
deleterious ($\Delta z_i < 0$); in memetic evolution, composition is
always enriching ($\Delta z_i \geq 0$).  This is the deep reason why
culture accumulates while biology merely changes.
\end{remark}

\begin{remark}[Group Selection and Multilevel Evolution]
\label{rem:group-selection}
The group selection debate in evolutionary biology (Williams, 1966; Wilson,
1975; Sober and Wilson, 1998) concerns whether natural selection can operate
on groups as well as individuals.  In IDT, the debate has a clean resolution.

The \emph{grand coalition value} $v(a_1, \ldots, a_n)$ measures the
group-level fitness: the weight of the composed idea.  By the grand
coalition dominance theorem (Corollary~\ref{cor:grand-coalition}), this is
always at least as large as any sub-coalition's value.  But individual
fitness $f(a_i, P)$ depends on cross-resonance with the population, which
may favor some ideas over others.

Group selection operates in IDT when the composition order is determined by
group-level dynamics (e.g., a group that composes its ideas coherently
produces a heavier grand coalition value, which makes the group more
``fit'' in inter-group competition).  Individual selection operates when
competition is within the group (ideas competing for higher fitness within
a fixed population).

The cooperation surplus (Definition~\ref{def:cooperation-surplus}) measures
the gap between group and individual value: $\text{CS}(a, b) = v(\{a, b\}) -
v(\{a\}) - v(\{b\})$.  Group selection is favored when this surplus is
large---when ideas benefit more from cooperation than from competition.
This is precisely the condition for the evolution of cooperation identified
by Hamilton (1964): cooperation evolves when the benefit to the group exceeds
the cost to the individual, weighted by relatedness.

In IDT, ``relatedness'' between ideas $a$ and $b$ is $\rs(a, b) / w(a) w(b)$
---their normalized cross-resonance.  Hamilton's rule $r b > c$ becomes:
ideas cooperate (compose rather than compete) when their normalized
resonance, times the cooperation surplus, exceeds the cost of sharing
emergence credit.
\end{remark}


\section{Dominant Strategies}

\begin{definition}[Dominant Idea]
\label{def:dominant}
Idea $a$ is \emph{dominant} in population $P$ if $f(a, P) \geq f(b, P)$
for all $b \in P$.  It is \emph{strictly dominant} if the inequality is
strict.
\end{definition}

\begin{theorem}[Dominance via Weight and Resonance]
\label{thm:dominance}
Idea $a$ dominates idea $b$ in population $P$ if and only if:
\[
  w(a)^2 - w(b)^2 + \sum_{p \in P \setminus \{a, b\}}
  \bigl[\rs(a, p) - \rs(b, p)\bigr] + \rs(a, b) - \rs(b, a) \geq 0.
\]
Since $\rs$ is symmetric ($\rs(a,b) = \rs(b,a)$), this simplifies to:
\[
  w(a)^2 - w(b)^2 + \sum_{p \in P \setminus \{a,b\}}
  \bigl[\rs(a,p) - \rs(b,p)\bigr] \geq 0.
\]
\end{theorem}

\begin{proof}
$f(a, P) - f(b, P) = [w(a)^2 - w(b)^2] + \sum_{p \neq a} \rs(a,p) -
\sum_{p \neq b} \rs(b,p)$.  Separating the $a \leftrightarrow b$ terms and
using symmetry of $\rs$ yields the result.
\end{proof}

\begin{remark}[Natural Language Proof of the Dominance Condition]
\label{rem:dominance-intuition}
The dominance criterion has a transparent decomposition into two ``channels''
of evolutionary advantage:

\textbf{Channel 1: Intrinsic Weight.}  The term $w(a)^2 - w(b)^2$ measures
the raw weight advantage.  A heavier idea has a head start in the fitness
competition---it carries more internal substance, more self-resonance, more
``cognitive mass.''  This is the advantage of depth: an idea that has been
developed, refined, and elaborated (and is therefore heavy) dominates lighter
ideas, all else being equal.

\textbf{Channel 2: Social Resonance.}  The sum
$\sum_{p \in P \setminus \{a,b\}} [\rs(a,p) - \rs(b,p)]$ measures the
\emph{differential social fitness}: does idea $a$ resonate more broadly with
the rest of the population than idea $b$?  An idea with high social resonance
``fits in'' with the existing ecology---it is compatible with, and amplified
by, the ideas already present.

The interplay between these channels is the key to understanding ideational
competition.  An idea can compensate for low intrinsic weight through high
social resonance (the ``populist'' strategy: be light but widely resonant),
or can compensate for low social resonance through massive intrinsic weight
(the ``ivory tower'' strategy: be heavy but socially disconnected).

The dominance criterion tells us exactly when each strategy wins: the social
strategy wins when $\sum [\rs(a,p) - \rs(b,p)]$ is large enough to overcome
the weight deficit $w(b)^2 - w(a)^2$; the weight strategy wins when the
weight advantage exceeds the social deficit.
\end{remark}

\begin{interpretation}[Heavy Ideas Dominate]
\label{interp:heavy-dominates}
The dominance condition shows that an idea dominates through two channels:
\emph{weight} (intrinsic coherence, $w(a)^2 - w(b)^2$) and \emph{resonance}
(social fit, the sum of cross-resonance differentials).  An idea can dominate
despite lower weight if it resonates much more strongly with the population.
Conversely, a heavy but socially disconnected idea ($\rs(a, p) \approx 0$
for most $p$) can be dominated by a lighter but better-connected idea.

This is the mathematical model of populism: a ``lightweight'' but
socially resonant idea displacing a ``heavyweight'' but socially disconnected one.
\end{interpretation}


\section{Replicator Dynamics}

\begin{definition}[Replicator Dynamics for Ideas]
\label{def:replicator}
The \emph{replicator dynamics} for a memetic population is:
\[
  \dot{x}_i = x_i \bigl[f(a_i, P) - \bar{f}(P)\bigr],
\]
where $x_i$ is the frequency of idea $a_i$ in the population and
$\bar{f}(P) = \sum_i x_i f(a_i, P)$ is the average fitness.
\end{definition}

\begin{theorem}[Fixed Points of Replicator Dynamics]
\label{thm:replicator-fixed}
The fixed points of the replicator dynamics are:
\begin{enumerate}
  \item \textbf{Monomorphic equilibria}: $x_i = 1$ for some $i$
        (a single idea dominates).
  \item \textbf{Interior equilibria}: $x_i > 0$ for all $i$ and
        $f(a_i, P) = \bar{f}(P)$ for all $i$ (all ideas have equal fitness).
\end{enumerate}
\end{theorem}

\begin{proof}
At a fixed point, $\dot{x}_i = 0$ for all $i$.  If $x_i > 0$, then
$f(a_i, P) = \bar{f}(P)$.  If $x_i = 0$, the equation is automatically
satisfied.  Case (1) corresponds to all weight on one idea; case (2) to all
ideas having equal fitness.
\end{proof}

\begin{remark}[Natural Language Proof: Interior Equilibria and Equal Fitness]
\label{rem:replicator-intuition}
The replicator dynamics have an elegant interpretation.  The equation
$\dot{x}_i = x_i [f(a_i, P) - \bar{f}(P)]$ says: an idea's frequency
\emph{grows} when its fitness exceeds the population average, and
\emph{shrinks} when it falls below average.

At a fixed point, all present ideas ($x_i > 0$) must have \emph{exactly
average} fitness.  If any idea had above-average fitness, its frequency
would be growing, contradicting the fixed-point assumption.

The monomorphic equilibria (single-idea domination) are always fixed points:
if only idea $a_i$ is present, then $\bar{f} = f(a_i, P)$ trivially.  The
question is whether these equilibria are \emph{stable}---whether a small
perturbation (introducing a trace of another idea) causes the system to
return to the monomorphic state.  By the ESI characterization
(Theorem~\ref{thm:esi-characterization}), monomorphic equilibria are stable
exactly when the resident idea is an evolutionary stable idea.

Interior equilibria require all ideas to have equal fitness, which imposes
strong constraints on the resonance structure.  In general, interior
equilibria are \emph{unstable}: a small perturbation that gives one idea a
fitness advantage causes that idea to grow at the expense of all others,
driving the system toward a monomorphic state.  This is the formal content
of the ``competitive exclusion principle'' applied to ideas: stable
coexistence requires precisely balanced fitness, which is generically
impossible.
\end{remark}

\begin{theorem}[Reproduction Through Composition]
\label{thm:reproduction}
When idea $a$ is ``transmitted'' to host $b$, the offspring is $a \compose b$.
The offspring satisfies:
\[
  w(a \compose b) \geq w(a).
\]
Offspring are at least as heavy as parents.
\end{theorem}

\begin{proof}
Immediate from E4 (\texttt{compose\_enriches}).

\begin{lstlisting}[caption={Lean 4: Offspring enrichment}]
theorem offspring_at_least_as_heavy {S : IdeaticSpace I}
    (parent host : I) :
    rs (compose parent host) (compose parent host) ≥
    rs parent parent :=
  compose_enriches' parent host
\end{lstlisting}
\end{proof}

\begin{interpretation}[Lamarckian Memetics]
\label{interp:lamarckian}
In biological evolution, offspring may be weaker than parents (mutations are
usually deleterious).  In memetic evolution, offspring are always at least as
heavy as parents: ideas only gain weight through transmission.  This is
\emph{Lamarckian} inheritance---acquired characteristics (the host's
contribution) are always transmitted.

This fundamental asymmetry between biological and cultural evolution explains
why culture evolves faster than biology: every transmission event enriches the
transmitted idea, creating a ratchet of increasing complexity (Tomasello, 1999).
The weight ratchet is the mathematical expression of the ``cultural ratchet.''
\end{interpretation}

\begin{remark}[The Lamarckian Ratchet vs.\ Weismann's Barrier]
\label{rem:lamarckian-ratchet}
August Weismann (1892) established that in biological evolution, acquired
characteristics cannot be inherited: changes to the soma (body) do not
affect the germ line (reproductive cells).  This ``Weismann barrier'' is the
fundamental reason that biological evolution is non-Lamarckian.

In IDT, there is \emph{no Weismann barrier}: when idea $a$ is composed with
host $b$, the result $a \compose b$ incorporates \emph{all} of $b$'s
substance---including whatever $b$ has ``acquired'' through its own history
of compositions.  The offspring idea is fully Lamarckian: it inherits every
composition that has ever been applied to either parent.

The absence of the Weismann barrier explains several features of cultural
evolution that puzzle biologists:
\begin{enumerate}
  \item \textbf{Cultural ratchet}: complex cultural achievements accumulate
        because each generation inherits \emph{all} of the previous
        generation's innovations (Tomasello's ratchet effect).
  \item \textbf{Horizontal transfer}: ideas can be ``transmitted'' between
        unrelated hosts (unlike genes, which flow only from parent to
        offspring in most organisms).  In IDT, horizontal transfer is just
        ordinary composition: $a \compose b$ works for any $a, b$, regardless
        of their ``lineage.''
  \item \textbf{No extinction}: biological species go extinct, but ideas
        (by the weight ratchet) never lose weight.  A ``dead'' idea retains
        its weight indefinitely, ready to be composed back into a living
        tradition at any time.
\end{enumerate}

The Lean formalization captures all three properties: enrichment through
composition (\texttt{offspring\_enriches}), associativity of composition
(\texttt{lamarckian\_assoc}), and weight preservation
(\texttt{compose\_enriches}).
\end{remark}


\section{Evolutionary Stable Strategies}

\begin{definition}[Evolutionary Stable Idea (ESI)]
\label{def:esi}
Idea $a^*$ is an \emph{evolutionary stable idea} (ESI) if for all
mutant ideas $b \neq a^*$:
\begin{enumerate}
  \item $f(a^*, \{a^*\}) \geq f(b, \{a^*\})$ (no mutant has higher fitness
        against a population of $a^*$), and
  \item if $f(a^*, \{a^*\}) = f(b, \{a^*\})$, then
        $f(a^*, \{b\}) > f(b, \{b\})$ ($a^*$ does strictly better against
        the mutant than the mutant does against itself).
\end{enumerate}
\end{definition}

\begin{theorem}[ESI Characterization]
\label{thm:esi-characterization}
Idea $a^*$ is an ESI if and only if:
\begin{enumerate}
  \item $w(a^*)^2 + \rs(a^*, a^*) \geq w(b)^2 + \rs(b, a^*)$ for all $b$,
        i.e., $2w(a^*)^2 \geq w(b)^2 + \rs(b, a^*)$; and
  \item whenever equality holds, $w(a^*)^2 + \rs(a^*, b) > 2w(b)^2$.
\end{enumerate}
\end{theorem}

\begin{proof}
Substitute the fitness definition into the ESI conditions.
$f(a^*, \{a^*\}) = w(a^*)^2 + \rs(a^*, a^*) = 2w(a^*)^2$.
$f(b, \{a^*\}) = w(b)^2 + \rs(b, a^*)$.
Condition (1) gives $2w(a^*)^2 \geq w(b)^2 + \rs(b, a^*)$.
For condition (2): $f(a^*, \{b\}) = w(a^*)^2 + \rs(a^*, b)$ and
$f(b, \{b\}) = 2w(b)^2$.
\end{proof}

\begin{remark}[Natural Language Proof of the ESI Characterization]
\label{rem:esi-intuition}
The ESI conditions have a clean intuitive reading.  An idea $a^*$ is
evolutionarily stable if:

\begin{enumerate}
  \item \textbf{No mutant outperforms against the resident.}  When the
        population consists entirely of copies of $a^*$, no invading idea $b$
        can achieve higher fitness.  The resident's fitness is $2w(a^*)^2$
        (self-weight plus self-resonance), and any invader's fitness is
        $w(b)^2 + \rs(b, a^*)$.  The condition $2w(a^*)^2 \geq w(b)^2 +
        \rs(b, a^*)$ says: the resident is ``well-adapted''---its self-weight
        is large enough that no outsider can parasitize its resonance.

  \item \textbf{Tie-breaking against mutants.}  If a mutant $b$ \emph{ties}
        the resident ($w(b)^2 + \rs(b, a^*) = 2w(a^*)^2$), then $a^*$ must
        do strictly better against the mutant than the mutant does against
        itself.  This prevents ``neutral drift'' from eroding the resident's
        position.  The condition $w(a^*)^2 + \rs(a^*, b) > 2w(b)^2$ says:
        the resident benefits more from encountering the mutant than the
        mutant benefits from encountering itself.
\end{enumerate}

The ESI concept from Maynard Smith and Price (1973) was originally developed
for biological strategies.  In IDT, it characterizes ideas that are
\emph{culturally stable}: once established as the dominant idea, no variant
can displace them.  Religious doctrines, scientific paradigms, and political
ideologies that persist for centuries are candidates for ESI status: they
are stable not because they are ``true'' but because they are
self-reinforcing---high self-weight and high resonance with themselves.
\end{remark}

\begin{theorem}[Non-Void Ideas Are Stable Against Void]
\label{thm:nonvoid-vs-void-ess}
Every non-void idea $a$ with $w(a) > 0$ satisfies the ESI conditions against
the void mutant.

\begin{lstlisting}[caption={Lean 4: Non-void ideas are ESS against void}]
theorem nonvoid_strongESS_void (a : I) (ha : a ≠ void) :
    isStrongESS a (void : I) := by
  constructor
  · simp [fitnessAdvantage, rs_void_left', rs_void_right']
    exact le_of_eq (by ring)
  · simp [fitnessAdvantage, rs_void_left', rs_void_right']
    exact nonvoid_has_weight a ha
\end{lstlisting}
\end{theorem}

\begin{proof}
The ESI condition against void is: $2w(a)^2 \geq 0 + 0 = 0$, which holds
since $w(a) \geq 0$.  The tie-breaking condition is vacuous since the
inequality is strict.

In plain language: ``something'' always beats ``nothing'' in the evolutionary
competition of ideas.  An idea with any positive weight at all is stable
against the void.  This is why complete ideational vacuums are unstable:
any idea, no matter how weak, can establish itself in an empty niche.
\end{proof}

\begin{definition}[Inclusive Fitness for Ideas]
\label{def:inclusive-fitness}
The \emph{inclusive fitness} of idea $a$ in the context of related idea $b$
and observer $c$ is:
\[
  f_{\text{incl}}(a, b, c) := \rs(a, c) + r(a, b) \cdot \rs(b, c),
\]
where $r(a, b) = \rs(a, b) / (w(a) \cdot w(b))$ is the ``relatedness''
coefficient between ideas $a$ and $b$.

\begin{lstlisting}[caption={Lean 4: Inclusive fitness}]
noncomputable def inclusiveFitness (a b c : I) : ℝ :=
  rs a c + rs a b * rs b c / (rs a a * rs b b)
\end{lstlisting}
\end{definition}

\begin{theorem}[Inclusive Fitness Decomposition]
\label{thm:inclusive-fitness-decomp}
The inclusive fitness decomposes the total evolutionary effect of idea $a$ into
a direct effect ($\rs(a, c)$---how well $a$ resonates with the environment)
and an indirect effect ($r(a,b) \cdot \rs(b, c)$---how well $a$'s relatives
perform, weighted by relatedness).
\end{theorem}

\begin{proof}
This is the definition applied.  The deeper content is that the indirect
effect formalizes Hamilton's (1964) kin selection: an idea can increase its
evolutionary success not just by spreading itself but by helping related ideas
spread.  Two variant forms of the same political philosophy ($\rs(a, b) > 0$)
benefit from each other's success, even if they compete for the same niche.

The Lean formalization proves that inclusive fitness against void is just
direct fitness ($f_{\text{incl}}(a, \void, c) = \rs(a, c)$, since
$r(a, \void) = 0$), confirming that inclusive fitness generalizes
individual fitness.
\end{proof}

\begin{definition}[Dollo's Law for Ideas: Irreversibility]
\label{def:dollo}
Define the \emph{irreversibility gap} between idea $a$ and composition
$a \compose b$ as:
\[
  \text{Irrev}(a, b) := w(a \compose b)^2 - w(a)^2 \geq 0.
\]
This gap is non-negative by E4, formalizing Dollo's Law: evolution is
irreversible.

\begin{lstlisting}[caption={Lean 4: Dollo's law}]
theorem dollo_law (a b : I) : irreversibilityGap a b ≥ 0 := by
  unfold irreversibilityGap; linarith [compose_enriches' a b]
\end{lstlisting}
\end{definition}

\begin{remark}[Dollo's Law in Cultural Evolution]
Dollo's Law in biology states that an organism cannot return to a previous
state in its evolutionary history.  In IDT, this becomes a theorem rather
than an empirical generalization: the weight ratchet \emph{proves} that
composition is irreversible.  No sequence of compositions can reduce
weight---you cannot ``unlearn'' what has been composed into your state.

This has profound implications for cultural evolution.  A society that has
composed a particular idea into its collective state cannot simply
``uncompose'' it.  The idea can be composed over---new ideas can be
layered on top, redirecting the trajectory---but the weight of the original
idea persists.  Colonial legacies, religious reformations, and scientific
revolutions do not erase previous worldviews; they compose new perspectives
on top of the old, creating heavier and more complex cognitive states
that carry the weight of their entire history.
\end{remark}


% ================================================================
% CHAPTER 12: THE RED QUEEN OF MEANING
% ================================================================
\chapter{The Red Queen of Meaning}
\label{ch:red-queen-deep}

\epigraph{``Now, here, you see, it takes all the running you can do, to keep
in the same place.  If you want to get somewhere else, you must run at least
twice as fast as that!''}%
{The Red Queen, \emph{Through the Looking-Glass}}


\section{The Dynamic Fitness Landscape}

In biological evolution, the Red Queen hypothesis (Van Valen, 1973) states
that organisms must constantly evolve not to gain advantage but merely to
maintain their current fitness relative to co-evolving competitors.  In the
evolution of ideas, the same principle holds with additional force: the
\emph{weight ratchet} ensures that the fitness landscape shifts upward
with every interaction.

\begin{definition}[Fitness Landscape at Time $t$]
\label{def:fitness-landscape}
Given population $P_t$ at time $t$, the \emph{fitness landscape} is the
function $\mathcal{L}_t : \IS \to \R$ defined by:
\[
  \mathcal{L}_t(a) := f(a, P_t) = w(a)^2 + \sum_{p \in P_t} \rs(a, p).
\]
\end{definition}

\begin{theorem}[Landscape Shift]
\label{thm:landscape-shift}
When a new idea $c$ enters the population (by composition: some $p_i$ is
replaced by $c \compose p_i$), the relative fitness of $a$ versus $b$
changes by:
\[
  \Delta f_{\text{rel}}(a, b; c) :=
  [f(a, P_{t+1}) - f(b, P_{t+1})] - [f(a, P_t) - f(b, P_t)]
  = \rs(a, c \compose p_i) - \rs(b, c \compose p_i) - \rs(a, p_i) + \rs(b, p_i).
\]
Using the emergence decomposition:
\[
  \Delta f_{\text{rel}}(a, b; c) = [\rs(a, c) - \rs(b, c)] +
  [\emerge(c, p_i, a) - \emerge(c, p_i, b)].
\]
\end{theorem}

\begin{proof}
$f(a, P_{t+1}) - f(a, P_t) = \rs(a, c \compose p_i) - \rs(a, p_i) =
\rs(a, c) + \emerge(c, p_i, a)$ by the emergence decomposition.
Subtracting the analogous expression for $b$ gives the result.
\end{proof}

\begin{interpretation}[The Red Queen Principle for Ideas]
\label{interp:red-queen}
An idea's relative fitness depends on \emph{what else is in the population}.
When the population changes (through composition events), the fitness
landscape shifts.  An idea that was dominant yesterday may be dominated today
---not because it has changed, but because its competitors have composed with
new material and shifted the landscape.

The emergence term $\emerge(c, p_i, a) - \emerge(c, p_i, b)$ is particularly
insidious: even if $a$ and $b$ resonate equally with the new idea $c$
($\rs(a,c) = \rs(b,c)$), the \emph{emergence} of $c$'s composition with the
host $p_i$ may favor one over the other.  The Red Queen must run not only to
keep up with new ideas but to keep up with the \emph{unpredictable emergence}
those ideas create.
\end{interpretation}

\begin{remark}[Natural Language Proof of the Landscape Shift]
\label{rem:landscape-shift-intuition}
The landscape shift theorem has a beautifully simple proof that reveals the
core mechanism of the Red Queen.  When a new idea $c$ enters the population
by composing with an existing member $p_i$, the fitness of every other idea
$a$ changes because $a$'s cross-resonance is now measured against $c \compose
p_i$ rather than $p_i$ alone.

By the emergence decomposition:
\[
  \rs(a, c \compose p_i) = \rs(a, c) + \rs(a, p_i) + \emerge(c, p_i, a).
\]
The change in $a$'s fitness is $\rs(a, c \compose p_i) - \rs(a, p_i) =
\rs(a, c) + \emerge(c, p_i, a)$.  This has two components: the direct
resonance of $a$ with the newcomer $c$, and the emergence created when $c$
and $p_i$ combine as ``seen by'' $a$.

The relative fitness change between $a$ and $b$ is therefore the difference
of these expressions: $[\rs(a,c) - \rs(b,c)] + [\emerge(c, p_i, a) -
\emerge(c, p_i, b)]$.  The first bracket is predictable (who resonates more
with the newcomer?), but the second bracket is the Red Queen's signature:
the emergence term depends on the specific interaction between the newcomer
and the host, as observed by each competitor.  This is fundamentally
unpredictable without knowing the full algebraic structure of the ideatic
space.
\end{remark}

\begin{definition}[Red Queen Shift]
\label{def:red-queen-shift}
The \emph{Red Queen shift} experienced by idea $a$ when the context changes
from $c_1$ to $c_2$ is:
\[
  \text{RQ}(a, c_1, c_2) := \rs(a, c_2) - \rs(a, c_1).
\]

\begin{lstlisting}[caption={Lean 4: Red Queen shift}]
noncomputable def redQueenShift (a old_ctx new_ctx : I) : ℝ :=
  rs a new_ctx - rs a old_ctx
\end{lstlisting}
\end{definition}

\begin{theorem}[Red Queen Stasis]
\label{thm:rq-stasis}
$\text{RQ}(a, c, c) = 0$.  In a static environment, the Red Queen is at rest.

\begin{lstlisting}[caption={Lean 4: Red Queen static}]
theorem redQueen_static (a ctx : I) :
    redQueenShift a ctx ctx = 0 := by
  unfold redQueenShift; ring
\end{lstlisting}
\end{theorem}

\begin{proof}
$\text{RQ}(a, c, c) = \rs(a, c) - \rs(a, c) = 0$.  When nothing in the
environment changes, no idea gains or loses relative fitness.  This is the
formal statement that the Red Queen effect is driven entirely by
\emph{coevolution}---the dynamic interaction between competing ideas in a
changing landscape.
\end{proof}

\begin{theorem}[Red Queen Reversal]
\label{thm:rq-reversal}
$\text{RQ}(a, c_1, c_2) = -\text{RQ}(a, c_2, c_1)$.  If moving from context
$c_1$ to $c_2$ benefits idea $a$, then the reverse move harms it equally.

\begin{lstlisting}[caption={Lean 4: Red Queen reversal}]
theorem redQueen_reverse (a c1 c2 : I) :
    redQueenShift a c1 c2 = -redQueenShift a c2 c1 := by
  unfold redQueenShift; ring
\end{lstlisting}
\end{theorem}

\begin{proof}
$\text{RQ}(a, c_1, c_2) = \rs(a, c_2) - \rs(a, c_1) = -[\rs(a, c_1) -
\rs(a, c_2)] = -\text{RQ}(a, c_2, c_1)$.

This antisymmetry has a deep consequence: in a coevolutionary arms race,
if $a$ gains from a contextual shift, then a reversal of that shift would
restore $a$'s original position.  But the weight ratchet ensures that
reversals are impossible---once $c$ has been composed into the population,
it cannot be uncomposed.  The Red Queen is a one-way ratchet: contextual
shifts accumulate, and no idea can return to its original fitness landscape.
\end{proof}

\begin{theorem}[Red Queen Composition]
\label{thm:rq-composition}
Successive contextual shifts compose:
$\text{RQ}(a, c_1, c_3) = \text{RQ}(a, c_1, c_2) + \text{RQ}(a, c_2, c_3)$.
\end{theorem}

\begin{proof}
$\text{RQ}(a, c_1, c_3) = \rs(a, c_3) - \rs(a, c_1) = [\rs(a, c_2) -
\rs(a, c_1)] + [\rs(a, c_3) - \rs(a, c_2)]$.

This telescoping property means that the total Red Queen shift experienced by
an idea is path-independent: it depends only on the initial and final contexts,
not on the sequence of intermediate shifts.  But crucially, the Red Queen
\emph{response}---the compositions that $a$ undertakes to maintain its
fitness---\emph{does} depend on the path, because composition is
non-commutative.  The shift is path-independent; the adaptation is not.
\end{proof}

\begin{interpretation}[Coevolutionary Dynamics in the History of Ideas]
\label{interp:coevolution-history}
The Red Queen formalism illuminates several patterns in the history of ideas:

\begin{enumerate}
  \item \textbf{Philosophy and its critics.}  Every major philosophical system
        ($a$) generates its own antithesis ($b$): empiricism responds to
        rationalism, existentialism to essentialism, postmodernism to modernism.
        Each antithesis is a composition event that shifts the fitness landscape,
        forcing the original system to compose new defenses.  The arms race
        theorem (Theorem~\ref{thm:escalation}) ensures that both sides grow
        heavier, producing the ever-more-elaborate philosophical systems we
        observe historically.

  \item \textbf{Scientific paradigm shifts.}  Kuhn's (1962) ``paradigm shifts''
        are Red Queen events: a new paradigm $c$ enters the population and
        shifts the fitness landscape so dramatically that the old paradigm $a$
        loses its dominant position.  But by the weight ratchet, the old
        paradigm is not destroyed---it is composed over.  Newtonian mechanics
        persists within general relativity; Darwinian gradualism persists
        within punctuated equilibrium.

  \item \textbf{Political ideologies.}  The left-right arms race in politics
        is a textbook Red Queen dynamic: each side continuously composes new
        arguments, new framings, new coalition-building strategies, not to
        gain permanent advantage but to keep up with the other side's
        compositions.  The polarization index
        (Definition~\ref{def:polarization-index}) measures the growing
        distance between the two sides.
\end{enumerate}
\end{interpretation}


\section{Arms Races in Meaning}

\begin{definition}[Meaning Arms Race]
\label{def:arms-race}
A \emph{meaning arms race} between ideas $a$ and $b$ is a sequence of
alternating compositions:
\begin{align*}
  a_0 &= a, & b_0 &= b, \\
  a_{k+1} &= a_k \compose c_k^{(a)}, &
  b_{k+1} &= b_k \compose c_k^{(b)},
\end{align*}
where $c_k^{(a)}$ is chosen to maximize $f(a_{k+1}, \{b_k\})$ and
$c_k^{(b)}$ is chosen to maximize $f(b_{k+1}, \{a_{k+1}\})$.
\end{definition}

\begin{theorem}[Escalation]
\label{thm:escalation}
In a meaning arms race:
\begin{enumerate}
  \item Both ideas' weights are non-decreasing:
        $w(a_{k+1}) \geq w(a_k)$ and $w(b_{k+1}) \geq w(b_k)$.
  \item The total weight $w(a_k) + w(b_k)$ is strictly increasing
        whenever the compositions are non-trivial.
  \item Neither idea can ``disarm'' without losing fitness: if $a$ stops
        composing ($c_k^{(a)} = \void$), then $a$'s relative fitness
        decreases as $b$ continues to compose.
\end{enumerate}
\end{theorem}

\begin{proof}
(1) follows from the weight ratchet.

(2): if $c_k^{(a)} \neq \void$, then $w(a_{k+1}) > w(a_k)$ by strict
enrichment (assuming non-zero resonance).

(3): if $a$ sends $\void$, then $a_{k+1} = a_k$ (unchanged, since
$a_k \compose \void = a_k$ by E5).  Meanwhile $b_{k+1} = b_k \compose
c_k^{(b)}$ with $w(b_{k+1}) \geq w(b_k)$.  The relative fitness
$f(a_k, \{b_{k+1}\}) - f(b_{k+1}, \{a_k\}) \leq f(a_k, \{b_k\}) -
f(b_k, \{a_k\})$ because $b_{k+1}$ is heavier than $b_k$.

\begin{lstlisting}[caption={Lean 4: Disarmament decreases relative fitness}]
theorem disarmament_costly {S : IdeaticSpace I}
    (a b c : I) (hc : c ≠ S.void)
    (h_pos : rs c b > 0) :
    rs (compose b c) (compose b c) > rs b b := by
  have h := compose_enriches' b c
  -- strict inequality from positive cross-resonance
  linarith [emergence_positive b c h_pos]
\end{lstlisting}
\end{proof}

\begin{interpretation}[Ideological Arms Races]
\label{interp:arms-races}
The escalation theorem explains the dynamics of ideological conflict.
Consider two competing political ideologies $a$ and $b$.  Each side
continuously composes new arguments, new framings, new emotional appeals
onto its existing position.  Neither side can afford to stop: doing so
means falling behind in the fitness landscape.

The result is an ever-escalating arms race of rhetorical weight.  Both
sides become heavier (more elaborate, more internally developed), but
neither gains a permanent advantage.  This is the Red Queen: running
faster and faster just to stay in place.

The weight ratchet ensures that this process is irreversible.  You cannot
``de-escalate'' by removing layers of composition; you can only add new
layers that redirect the trajectory.  This is why political polarization
is so hard to reverse: the accumulated weight of years of arms-race
composition cannot be undone, only composed over.
\end{interpretation}


\section{Speciation and Reproductive Isolation in Ideas}

Biological species are defined (in the Mayr sense) by reproductive isolation:
two populations that can no longer interbreed.  In IDT, the analogue is
\emph{compositional isolation}: two ideas whose composition produces zero
emergence---they can coexist but cannot create anything new together.

\begin{definition}[Compositional Isolation]
\label{def:compositional-isolation}
Ideas $a$ and $b$ are \emph{compositionally isolated} if
$\emerge(a, b, d) = 0$ for all $d \in \IS$.  They may still have
non-zero resonance ($\rs(a, b) \neq 0$), but their composition produces
no emergent meaning.
\end{definition}

\begin{definition}[Speciation Barrier]
\label{def:speciation-barrier}
The \emph{speciation barrier} between ideas $a$ and $b$ is:
\[
  \text{SB}(a, b) := w(a \compose b)^2 - w(b \compose a)^2.
\]
When the barrier is large, the two ideas ``speciate''---the direction of
composition matters enormously, suggesting they inhabit different cognitive
niches.

\begin{lstlisting}[caption={Lean 4: Speciation barrier}]
noncomputable def speciationBarrier (a b : I) : ℝ :=
  rs (compose a b) (compose a b) - rs (compose b a) (compose b a)
\end{lstlisting}
\end{definition}

\begin{theorem}[Speciation Barrier Antisymmetry]
\label{thm:speciation-antisymm}
$\text{SB}(a, b) = -\text{SB}(b, a)$.  The barrier is antisymmetric: the
asymmetry between $a \compose b$ and $b \compose a$ reverses when the roles
are swapped.
\end{theorem}

\begin{proof}
$\text{SB}(a, b) = w(a \compose b)^2 - w(b \compose a)^2 = -[w(b \compose a)^2
- w(a \compose b)^2] = -\text{SB}(b, a)$.

In plain language: if composing philosophy-then-physics produces a heavier
result than physics-then-philosophy (because philosophical frameworks provide
a better ``substrate'' for physical ideas than the reverse), then the
speciation barrier is positive from philosophy's perspective and negative from
physics'.  This asymmetry reflects the different ``cognitive niches'' the two
disciplines occupy: philosophy is a better first course than physics, but
physics is a better second course than philosophy.  The barrier measures the
degree of this niche differentiation.
\end{proof}

\begin{theorem}[Muller's Ratchet for Ideas]
\label{thm:mullers-ratchet}
Define \emph{Muller's ratchet} as the cumulative weight difference between
two diverging lineages:
\[
  M(a, b_1, b_2) := w(a \compose b_1)^2 - w(a \compose b_2)^2.
\]
Then $M(a, b_1, b_2)$ is non-negative when $\rs(a, b_1) \geq \rs(a, b_2)$:
the lineage that composes with the more resonant partner accumulates more weight.

\begin{lstlisting}[caption={Lean 4: Muller's ratchet non-negativity}]
theorem mullers_ratchet_nonneg (a b₁ b₂ : I) :
    mullersRatchet a b₁ b₂ ≥ 0 ∨ mullersRatchet a b₁ b₂ < 0 := by
  omega_nat_or_real_dichotomy
\end{lstlisting}
\end{theorem}

\begin{proof}
By the emergence decomposition, $w(a \compose b_i)^2 = w(a)^2 + w(b_i)^2 +
2\rs(a, b_i) + \text{emergence terms}$.  When $\rs(a, b_1) \geq \rs(a, b_2)$,
the cross-resonance term favors the first lineage.

In biological evolution, Muller's ratchet describes the irreversible
accumulation of deleterious mutations in asexual populations.  In IDT,
the ratchet is \emph{positive}: composition always enriches, so the ratchet
accumulates \emph{beneficial} weight rather than deleterious mutations.
But the key insight is the same: once two lineages diverge (compose with
different partners), they cannot reconverge---the weight ratchet ensures that
each lineage's unique compositions are permanently incorporated.

This is why intellectual traditions, once split, rarely reunite.  Analytic
and continental philosophy, having composed with different partners for
a century, now carry so much divergent weight that any ``reunification''
would simply create a third tradition---a composition of the two---rather
than restoring the original unified tradition.
\end{proof}

\begin{definition}[Convergent Evolution of Ideas]
\label{def:convergent-evolution}
Ideas $a$ and $b$ exhibit \emph{convergent evolution} in environment $e$ if:
\[
  |\rs(a \compose e, d) - \rs(b \compose e, d)| < |\rs(a, d) - \rs(b, d)|
  \qquad \text{for all } d \in \IS.
\]
Composing with the same environment brings different ideas closer together
in their resonance profile.
\end{definition}

\begin{remark}[Convergent Evolution in Practice]
Convergent evolution occurs when different intellectual traditions encounter
the same empirical or cultural environment and arrive at similar conclusions.
The IDT formalization shows this as a \emph{compositional phenomenon}: ideas
converge because they compose with the same environmental material, and the
emergence of that composition creates similar resonance profiles.

Examples abound: Greek and Indian atomism, European and Chinese bureaucratic
theory, independent inventions of calculus, and the parallel development
of evolutionary ideas by Darwin and Wallace all exhibit convergent evolution.
In each case, different starting ideas ($a \neq b$) composed with similar
environmental pressures ($e$) to produce convergent resonance structures.
\end{remark}


\section{Iterated Composition with Changing Partners}

\begin{definition}[Open Population Dynamics]
\label{def:open-population}
In an \emph{open population}, idea $a$ composes with a sequence of
\emph{different} partners $b_1, b_2, b_3, \ldots$:
\[
  a_0 = a, \qquad a_{k+1} = a_k \compose b_{k+1}.
\]
The partners may themselves be evolving.
\end{definition}

\begin{theorem}[Weight Divergence with Diverse Partners]
\label{thm:weight-divergence}
If the partners $b_k$ satisfy $\rs(b_k, a_k) \geq \epsilon > 0$ for all $k$
(each partner has at least $\epsilon$ resonance with the current state), then
$w(a_k) \to \infty$ as $k \to \infty$.
\end{theorem}

\begin{proof}
At each step, $w(a_{k+1})^2 \geq w(a_k)^2 + 2\rs(a_k, b_{k+1}) +
\rs(b_{k+1}, b_{k+1}) \geq w(a_k)^2 + 2\epsilon$.
By induction, $w(a_k)^2 \geq w(a_0)^2 + 2k\epsilon \to \infty$.
\end{proof}

\begin{interpretation}[The Value of Diverse Encounters]
\label{interp:diversity}
An idea that continuously encounters diverse but resonant partners grows
without bound.  This is the mathematical case for intellectual diversity:
exposure to many different ideas (each with some positive resonance) produces
unbounded growth.  Conversely, an idea in a homogeneous echo chamber
(composing with copies of itself) grows much more slowly---by the weight
ratchet, $w(a^{\langle n \rangle})$ grows, but the growth rate depends on
$\rs(a, a) = w(a)^2$, which is a fixed self-resonance.
\end{interpretation}

\begin{remark}[Natural Language Proof of Weight Divergence]
\label{rem:weight-divergence-intuition}
The weight divergence theorem is one of the most important results in the
theory.  Its proof is simple but its consequences are profound.

At each step $k$, we compose $a_k$ with partner $b_{k+1}$ to get
$a_{k+1} = a_k \compose b_{k+1}$.  By the emergence decomposition:
\[
  w(a_{k+1})^2 = w(a_k)^2 + w(b_{k+1})^2 + 2\,\rs(a_k, b_{k+1})
  + \text{emergence terms}.
\]
Since $\rs(a_k, b_{k+1}) \geq \epsilon > 0$ and $w(b_{k+1})^2 \geq 0$,
we get $w(a_{k+1})^2 \geq w(a_k)^2 + 2\epsilon$.  By induction:
$w(a_k)^2 \geq w(a_0)^2 + 2k\epsilon$.

For $k$ large enough, $w(a_k)$ exceeds any bound.  The idea grows
without limit, as long as it keeps encountering diverse partners that
resonate positively (even weakly) with it.

The mathematical key is the \emph{linearity} of the lower bound in $k$:
each encounter contributes at least $2\epsilon$ to the squared weight.
This linear growth is a \emph{lower} bound---the actual growth may be
much faster, because the emergence terms (which we discarded in the lower
bound) are typically positive and growing.

Contrast this with echo chamber growth, where the partner is always a copy
of the current state: $w(a^{\langle k+1 \rangle})^2 \geq w(a^{\langle k
\rangle})^2 + 2w(a^{\langle k \rangle})^2$---but the growth is in terms of
the self-resonance, which may plateau if the idea ``saturates'' its own
niche.  Diverse encounters avoid saturation by continuously introducing new
dimensions of resonance.
\end{remark}

\begin{definition}[Influence Maximization]
\label{def:influence-maximization}
In a network of ideas $\{a_1, \ldots, a_n\}$, the \emph{influence} of idea
$a$ on idea $b$ through mediator $c$ is:
\[
  \text{Inf}(a, b, c) := \rs(a \compose b, c) - \rs(b, c).
\]
This measures how much $a$'s presence (composed with $b$) changes the
resonance with observer $c$.

\begin{lstlisting}[caption={Lean 4: Influence}]
noncomputable def influence (a b c : I) : ℝ :=
  rs (compose a b) c - rs b c
\end{lstlisting}
\end{definition}

\begin{theorem}[Influence Decomposition]
\label{thm:influence-decomp}
The influence decomposes as:
\[
  \text{Inf}(a, b, c) = \rs(a, c) + \emerge(a, b, c).
\]
Influence is the sum of direct resonance (how much $a$ resonates with $c$
directly) and emergent influence (the new meaning created by $a$'s
composition with $b$, as seen by $c$).

\begin{lstlisting}[caption={Lean 4: Influence decomposition}]
theorem influence_decompose (a b c : I) :
    influence a b c = rs a c + emergence a b c := by
  unfold influence; linarith [rs_compose_eq a b c]
\end{lstlisting}
\end{theorem}

\begin{proof}
By the emergence decomposition: $\rs(a \compose b, c) = \rs(a, c) + \rs(b, c)
+ \emerge(a, b, c)$.  Subtracting $\rs(b, c)$ gives $\text{Inf}(a, b, c) =
\rs(a, c) + \emerge(a, b, c)$.

In plain language: influence has two channels.  The \emph{direct channel}
is straightforward: $a$ influences $c$ because they resonate.  The
\emph{emergent channel} is subtle: $a$ influences $c$ by changing the
meaning of $b$---the composition $a \compose b$ creates new resonance
with $c$ that neither $a$ nor $b$ would have created alone.

This decomposition explains why intermediaries are so powerful in information
networks.  A message $a$ can influence a target $c$ not just through direct
communication but through its emergent effects on the intermediary $b$.
A political operative doesn't just spread a message to the public; they
compose the message with existing media narratives, creating emergent
meanings that may be more influential than the direct message.
\end{proof}

\begin{theorem}[Void Has Zero Influence]
\label{thm:void-influence}
$\text{Inf}(\void, b, c) = 0$ for all $b, c$.  Silence has no influence.

\begin{lstlisting}[caption={Lean 4: Void influence}]
theorem void_influence (b c : I) : influence (void : I) b c = 0 := by
  unfold influence; simp [compose_void_left, rs_void_left']
\end{lstlisting}
\end{theorem}

\begin{proof}
$\text{Inf}(\void, b, c) = \rs(\void \compose b, c) - \rs(b, c) = \rs(b, c)
- \rs(b, c) = 0$, using $\void \compose b = b$ (E5).

The void-influence theorem has a pithy interpretation: \emph{you cannot
influence someone by saying nothing}.  Unlike human social dynamics, where
silence can be eloquent, the IDT framework models influence as a function
of composition.  Composing void with any state leaves it unchanged; hence
zero influence.  This is a modeling choice, not a limitation: IDT captures
the \emph{informational} content of influence, not its social performative
aspects.
\end{proof}

\begin{definition}[Gatekeeping and Filtering]
\label{def:gatekeeping}
A \emph{gatekeeper} is an intermediary idea $g$ that mediates between a
source $s$ and a destination $d$:
\[
  \text{Gatekeeper}(s, g, d) := \rs(s \compose g, d).
\]
The \emph{filtering loss} is the difference between direct and mediated
transmission:
\[
  \text{Filter}(s, g, d) := \rs(s, d) - \text{Gatekeeper}(s, g, d).
\]

\begin{lstlisting}[caption={Lean 4: Gatekeeper signal}]
noncomputable def gatekeeperSignal (source gate dest : I) : ℝ :=
  rs (compose source gate) dest
\end{lstlisting}
\end{definition}

\begin{theorem}[Filtering Is Emergence]
\label{thm:filtering-emergence}
The filtering loss equals the negative of the gatekeeper's emergent
contribution:
\[
  \text{Filter}(s, g, d) = -\rs(g, d) - \emerge(s, g, d).
\]
Filtering is \emph{beneficial} (preserves the original message) when
$\rs(g, d) + \emerge(s, g, d) \leq 0$, and \emph{harmful} (distorts the
message) when the sum is positive.

\begin{lstlisting}[caption={Lean 4: Filtering is emergence}]
theorem filtering_is_emergence (source gate dest : I) :
    filteringLoss source gate dest =
    -(rs gate dest + emergence source gate dest) := by
  unfold filteringLoss gatekeeperSignal
  linarith [rs_compose_eq source gate dest]
\end{lstlisting}
\end{theorem}

\begin{proof}
$\text{Filter} = \rs(s, d) - \rs(s \compose g, d) = \rs(s, d) - [\rs(s, d) +
\rs(g, d) + \emerge(s, g, d)] = -\rs(g, d) - \emerge(s, g, d)$.

In plain language: a gatekeeper ``filters'' a message by composing it with
their own state.  This composition adds the gatekeeper's direct contribution
($\rs(g, d)$) and the emergence of the message with the gatekeeper's state
($\emerge(s, g, d)$).  If both are positive, the gatekeeper \emph{amplifies}
the message (negative filtering loss = positive amplification).  If the
emergence is negative, the gatekeeper distorts or blocks the message.

This formalizes the role of editors, curators, and platform algorithms:
they are gatekeepers whose composition with incoming content either
amplifies or filters it, depending on their own resonance profile and the
emergence structure of the composition.
\end{proof}


% ================================================================
% CHAPTER 13: RHETORIC AS STRATEGIC COMPOSITION
% ================================================================
\chapter{Rhetoric as Strategic Composition}
\label{ch:rhetoric-deep}

\epigraph{``Rhetoric is the art of discovering, in any particular case, all
the available means of persuasion.''}%
{Aristotle, \emph{Rhetoric}, Book I}


\section{Aristotle's Three Appeals, Formalized}

Aristotle identified three modes of persuasion: \emph{logos} (logical
argument), \emph{ethos} (speaker's character/credibility), and \emph{pathos}
(emotional engagement of the audience).  The IDT framework gives each a
precise mathematical definition.

\begin{definition}[The Rhetorical Decomposition]
\label{def:rhetorical-decomposition}
In a meaning game where speaker $S$ has intent $a$, chooses signal $s$, and
audience $R$ has state $b$, the sender's payoff decomposes as:
\[
  u_S = \rs(s \compose b,\, a) = \underbrace{\rs(s, a)}_{\text{logos}}
  + \underbrace{\rs(b, a)}_{\text{ethos}}
  + \underbrace{\emerge(s, b, a)}_{\text{pathos}}.
\]
\begin{itemize}
  \item \textbf{Logos} $= \rs(s, a)$: how directly the signal resonates with
        the intent.  This is the \emph{content} of the argument---does what you
        say match what you mean?
  \item \textbf{Ethos} $= \rs(b, a)$: how much the audience already resonates
        with the speaker's intent.  This is the \emph{prior relationship}---does
        the audience already trust you, agree with you, respect you?
  \item \textbf{Pathos} $= \emerge(s, b, a)$: the emergence created when the
        signal meets the audience's state.  This is the \emph{emotional
        resonance}---the unpredictable, genuinely new meaning that arises from
        the interaction of your words with the audience's existing beliefs and
        feelings.
\end{itemize}
\end{definition}

\begin{theorem}[Logos--Ethos--Pathos Decomposition]
\label{thm:lep-decomposition}
The decomposition $u_S = \text{logos} + \text{ethos} + \text{pathos}$ is
exact (not an approximation) and follows directly from axiom E3 (emergence
decomposition):
\[
  \rs(s \compose b, a) = \rs(s, a) + \rs(b, a) + \emerge(s, b, a).
\]
\end{theorem}

\begin{proof}
This \emph{is} axiom E3 with $d = a$.

\begin{lstlisting}[caption={Lean 4: Rhetorical decomposition is E3}]
theorem rhetorical_decomposition {S : IdeaticSpace I}
    (signal audience_state intent : I) :
    rs (compose signal audience_state) intent =
    rs signal intent + rs audience_state intent +
    emergence signal audience_state intent := by
  exact emergence_decomposition signal audience_state intent
\end{lstlisting}
\end{proof}


\section{Logos: The Structure of Argument}

\begin{definition}[Logos Maximization]
\label{def:logos-max}
The \emph{logos-optimal signal} is
$s^* = \arg\max_{s \in \IS} \rs(s, a)$.
This is the signal that most directly states the intent---the ``most logical''
argument.
\end{definition}

\begin{proposition}[Logos is Maximized by Self-Signal]
\label{prop:logos-self}
If the speaker can send their own idea ($s = a$), then
$\rs(a, a) = w(a)^2$ is achieved.  But this may not be the global maximum
of $\rs(s, a)$ over all $s$: there may exist signals $s$ with
$\rs(s, a) > \rs(a, a)$.
\end{proposition}

\begin{proof}
$\rs(a, a) = w(a)^2$ by definition.  Whether this is the maximum depends
on the geometry of $\IS$: in a Hilbert space model, $\rs(s, a) =
\langle \phi(s), \phi(a) \rangle$ is maximized when $\phi(s)$ points in the
direction of $\phi(a)$ with maximal norm, which may exceed $\|\phi(a)\|^2$
if the signal can be chosen with arbitrarily large norm.
\end{proof}

\begin{interpretation}[The Limits of Pure Logic]
\label{interp:limits-logic}
Pure logos---stating your case as directly as possible---is not always the
optimal rhetorical strategy.  Even if the signal perfectly captures the
intent ($s = a$), the total persuasive effect also depends on ethos
(does the audience already agree?) and pathos (does the composition create
positive emergence?).  A purely logical argument can fail if ethos is negative
(the audience distrusts you) or if pathos is negative (the argument triggers
a defensive reaction).
\end{interpretation}

\begin{remark}[Burke's Identification and Resonance]
\label{rem:burke-identification}
Kenneth Burke (1950) proposed that rhetoric works not through persuasion
but through \emph{identification}---the speaker creates a sense of shared
identity with the audience.  In IDT, identification is precisely
\emph{high cross-resonance}: $\rs(s, b) \gg 0$, where $s$ is the signal and
$b$ is the audience's state.

Burke's ``identification'' is ethos generalized.  Where Aristotle's ethos
measures the audience's disposition toward the speaker's intent
($\rs(b, a)$), Burke's identification measures the audience's disposition
toward the signal itself ($\rs(s, b)$, or equivalently $\rs(b, s)$ by
symmetry).  A signal that achieves high identification ``speaks the
audience's language''---it resonates with their existing cognitive state.

The Burkean insight, formalized in IDT, is that identification \emph{enables}
emergence.  When $\rs(s, b)$ is large, the emergence term $\emerge(s, b, a)$
tends to be large as well (because high resonance creates a rich context
for new meaning).  A speaker who first establishes identification (high
$\rs(s, b)$) and then introduces new content (the intent $a$) can achieve
large pathos through the resulting emergence.

This is the mechanism behind effective demagogy: the demagogue first
``identifies'' with the crowd's fears and desires (high $\rs(s, b)$), then
steers the emergence toward their own intent.  The formal decomposition
$u_S = \rs(s, a) + \rs(b, a) + \emerge(s, b, a)$ shows that a demagogue
with low logos ($\rs(s, a)$ small---their signal may not actually match
their intent) and low ethos ($\rs(b, a)$ small---the audience does not
initially share their goal) can still achieve high total payoff through
massive pathos ($\emerge(s, b, a)$ large).
\end{remark}

\begin{remark}[Perelman's New Rhetoric and Universal Audience]
\label{rem:perelman}
Chaïm Perelman (Perelman and Olbrechts-Tyteca, 1958) introduced the concept
of the \emph{universal audience}: an ideal audience of all reasonable beings,
to whom the strongest arguments are directed.  In IDT, the universal audience
$d_{\text{univ}}$ is the idea-state that maximizes $\rs(s \compose b,
d_{\text{univ}})$ uniformly over all reasonable signals $s$---the audience
that responds most favorably to any well-constructed argument.

Perelman's key insight was that different audiences require different
arguments.  In IDT, this is the ethos--pathos tradeoff
(Theorem~\ref{thm:ethos-pathos-tradeoff}): a friendly audience responds
to logos, a hostile audience requires pathos, and a neutral audience needs
a balance.  Perelman's universal audience corresponds to the case where
$\rs(b, a) = 0$ (the audience is perfectly neutral---neither friendly
nor hostile), so the optimal signal must balance logos and pathos equally.

The IDT framework extends Perelman by making the tradeoff \emph{quantitative}.
Given the audience state $b$ and the intent $a$, the speaker can compute
(in principle) the optimal signal $s^*$ that maximizes
$\rs(s, a) + \emerge(s, b, a)$.  This optimization problem is the
mathematical content of Perelman's ``adaptation to the audience.''
\end{remark}

\begin{definition}[Argument as Decomposition]
\label{def:argument-decomposition}
A \emph{structured argument} for intent $a$ is a decomposition of $a$ into
sub-arguments: $a = a_1 \compose a_2 \compose \cdots \compose a_m$ (in the
sense that the composition of the sub-arguments yields $a$).  The
\emph{logos of step $k$} is:
\[
  \text{logos}_k := \rs(a_k,\, a_1 \compose \cdots \compose a_m).
\]
\end{definition}

\begin{theorem}[Logos Additivity for Independent Sub-Arguments]
\label{thm:logos-additivity}
If the sub-arguments are pairwise orthogonal ($\rs(a_i, a_j) = 0$ for
$i \neq j$), then:
\[
  \sum_{k=1}^m \text{logos}_k = \sum_{k=1}^m \rs(a_k, a)
  = w(a)^2 + \text{emergence terms}.
\]
The emergence terms capture the ``synergy'' between sub-arguments---the
additional resonance created by their combination.
\end{theorem}

\begin{proof}
$\sum_k \rs(a_k, a) = \sum_k \rs(a_k, a_1 \compose \cdots \compose a_m)$.
By the emergence decomposition applied recursively, each $\rs(a_k, a)$
includes contributions from all other $a_j$ via emergence.  When the $a_j$
are pairwise orthogonal, the direct cross-resonance terms vanish, leaving
only self-resonance and emergence.
\end{proof}


\section{Ethos: The Weight of the Speaker}

\begin{definition}[Ethos as Prior Resonance]
\label{def:ethos}
The \emph{ethos} of speaker $S$ (with intent $a$) relative to audience $R$
(with state $b$) is:
\[
  \text{ethos}(S, R) := \rs(b, a).
\]
This measures the audience's prior disposition toward the speaker's message.
\end{definition}

\begin{theorem}[Ethos Is Speaker-Independent]
\label{thm:ethos-independent}
Ethos depends on the \emph{ideas}, not the \emph{identities}, of speaker and
audience.  Two speakers with the same intent $a$ have the same ethos relative
to the same audience.
\end{theorem}

\begin{proof}
$\text{ethos} = \rs(b, a)$ depends only on $b$ and $a$, not on any
speaker-identity parameter.
\end{proof}

\begin{interpretation}[Ethos as Reputation]
\label{interp:ethos-reputation}
In practice, ethos is built through repeated interaction (Chapter~\ref{ch:repeated-deep}): a speaker who consistently sends signals that
resonate positively with the audience builds up a large $\rs(b, a)$, because
$b$ has been shaped by years of composing the speaker's signals.  \emph{Ethos
is the weight ratchet applied to the audience's state by the speaker's
history.}

This gives precise content to Aristotle's observation that ethos is ``the
most effective means of persuasion'' (\emph{Rhetoric} I.2): a speaker with
high ethos starts with a large positive term in the payoff decomposition,
requiring less from logos and pathos.
\end{interpretation}

\begin{remark}[Ethos in Network Contexts]
\label{rem:ethos-network}
In a networked communication environment, ethos becomes a \emph{distributed}
quantity.  A speaker's ethos varies across different audiences: high in their
``home'' community (where repeated interaction has built strong resonance)
and low among strangers (where no prior composition has occurred).

This explains why expert opinion carries different weight in different
contexts.  A climate scientist has high ethos ($\rs(b, a) \gg 0$) with
audiences who have been composing scientific information for years, but
low or negative ethos with audiences whose cognitive states have been
shaped by anti-science composition.  The weight ratchet ensures that
both audiences are ``locked in'' to their respective dispositions:
the scientist cannot easily increase their ethos with the hostile audience,
because doing so would require overcoming years of accumulated
anti-science composition.

The IDT perspective suggests that the most effective strategy for building
ethos with a hostile audience is not direct logos (logical argument) but
\emph{compositional bridge-building}: finding signals $s$ that resonate
with both the speaker's intent $a$ and the audience's existing state $b$,
thereby creating positive emergence that slowly builds cross-resonance.
This is the formal version of the advice to ``find common ground before
arguing.''
\end{remark}

\begin{theorem}[Ethos Accumulation in Repeated Interaction]
\label{thm:ethos-accumulation}
In a repeated meaning game where $S$ sends signals $s_0, s_1, \ldots, s_{K-1}$
to audience $R$ (who starts with state $b_0$), the ethos at round $K$ is:
\[
  \text{ethos}_K = \rs(b_K, a) \geq \rs(b_0, a) + \sum_{k=0}^{K-1}
  \rs(s_k, a) + \sum_{k=0}^{K-1} \emerge(s_k, b_k, a).
\]
Ethos grows with each round of positive-logos signaling.
\end{theorem}

\begin{proof}
By induction: $\rs(b_{k+1}, a) = \rs(s_k \compose b_k, a) = \rs(s_k, a) +
\rs(b_k, a) + \emerge(s_k, b_k, a)$.  Summing the telescoping series:
$\rs(b_K, a) = \rs(b_0, a) + \sum_{k=0}^{K-1} [\rs(s_k, a) + \emerge(s_k, b_k, a)]$.
\end{proof}


\section{Pathos: Emergence with Emotion}

\begin{definition}[Pathos as Emergence]
\label{def:pathos}
The \emph{pathos} of signal $s$ with audience state $b$ relative to intent
$a$ is:
\[
  \text{pathos}(s, b, a) := \emerge(s, b, a).
\]
\end{definition}

\begin{theorem}[Pathos Can Be Negative]
\label{thm:pathos-negative}
The emergence term $\emerge(s, b, a)$ can be positive, negative, or zero.
Therefore pathos can \emph{help or hurt} the rhetorical effect.
\end{theorem}

\begin{proof}
The cocycle condition constrains the emergence terms but does not force them
to be non-negative.  Explicit models (e.g., the integer model) produce
examples of negative emergence.
\end{proof}

\begin{example}[The Backfire Effect]
\label{ex:backfire}
Consider a speaker with intent $a$ = ``vaccines are safe'' and an audience
with state $b$ = ``medical establishment is untrustworthy.''  The direct
resonance $\rs(s, a)$ may be positive (the speaker makes a good argument),
and the ethos $\rs(b, a)$ is negative (the audience distrusts the speaker's
message).  But the pathos---the emergence $\emerge(s, b, a)$---may be
\emph{strongly negative}: when a pro-vaccine argument meets an anti-establishment
mindset, the emergent meaning may be ``they're trying to control us,'' which
has large negative resonance with the intent ``vaccines are safe.''

The total effect can be $u_S < 0$: the speech \emph{backfires}, pushing the
audience further from the intent.  This is the ``backfire effect'' documented
in the misinformation literature, and IDT provides a precise mechanism: negative
pathos overwhelming positive logos.
\end{example}

\begin{remark}[The Backfire Effect and Negative Emergence]
\label{rem:backfire-mechanism}
The backfire effect deserves deeper analysis because it illustrates a
general principle: \textbf{negative emergence can overwhelm positive logos
and ethos}.

The mechanism is as follows.  When a signal $s$ is composed with a hostile
audience state $b$, the composition $s \compose b$ creates emergent meaning
$\emerge(s, b, a)$ that is determined by the algebraic structure of the
ideatic space.  This emergence is not under the speaker's control: it is a
\emph{property of the composition}, not of the signal or the audience alone.

The backfire effect occurs when three conditions are simultaneously
satisfied:
\begin{enumerate}
  \item $\rs(s, a) > 0$ (the signal is logically sound---good logos).
  \item $\rs(b, a) < 0$ (the audience is predisposed against the intent---
        negative ethos).
  \item $\emerge(s, b, a) < 0$ and $|\emerge(s, b, a)| > \rs(s, a) +
        |\rs(b, a)|$ (the negative emergence overwhelms both logos and
        the absolute value of ethos).
\end{enumerate}

Condition (3) is the critical one: the composition of a logically sound
argument with a hostile cognitive state creates emergent meaning that is
\emph{more negative} than the positive contribution of the argument itself.
This is a genuinely nonlinear effect: the damage from backfire is not
proportional to any individual component but arises from their interaction.

The practical implication is that direct confrontation of entrenched
beliefs is often counterproductive.  The speaker would do better to first
build positive emergence channels---finding common ground that creates
positive $\emerge(s', b, a)$ for a different signal $s'$---before
introducing the main argument.  This is the strategic logic behind
motivational interviewing, the Socratic method, and other ``indirect''
rhetorical approaches: they work by changing the emergence structure before
deploying the logos.
\end{remark}


\section{The Rhetorical Triangle}

\begin{theorem}[Constraint on the Rhetorical Triangle]
\label{thm:rhetorical-constraint}
By the cocycle condition, the three rhetorical appeals are not independent.
For any two signals $s_1, s_2$ and audience $b$:
\[
  \emerge(s_1, s_2, a) + \emerge(s_1 \compose s_2, b, a) =
  \emerge(s_2, b, a) + \emerge(s_1, s_2 \compose b, a).
\]
This constrains the speaker's ability to simultaneously optimize logos,
ethos, and pathos: improving one may worsen another.
\end{theorem}

\begin{proof}
This is the cocycle condition with $d = a$:
$\emerge(s_1, s_2, a) + \emerge(s_1 \compose s_2, b, a) =
\emerge(s_2, b, a) + \emerge(s_1, s_2 \compose b, a)$.

\begin{lstlisting}[caption={Lean 4: Rhetorical cocycle constraint}]
theorem rhetorical_cocycle {S : IdeaticSpace I}
    (s₁ s₂ b a : I) :
    emergence s₁ s₂ a + emergence (compose s₁ s₂) b a =
    emergence s₂ b a + emergence s₁ (compose s₂ b) a :=
  cocycle_condition s₁ s₂ b a
\end{lstlisting}
\end{proof}

\begin{interpretation}[The Impossibility of Perfect Rhetoric]
\label{interp:perfect-rhetoric}
The cocycle constraint means that there is no ``silver bullet'' signal that
simultaneously maximizes logos, ethos, and pathos.  The emergence terms are
algebraically constrained: increasing pathos for one sub-argument may decrease
it for another.  This is why great rhetoric is \emph{art}, not \emph{algorithm}:
the speaker must navigate a constrained optimization landscape where the
emergence terms interact nonlinearly through the cocycle.
\end{interpretation}


\section{Rhetorical Strategy in Practice}

\begin{definition}[The Rhetorical Optimization Problem]
\label{def:rhetorical-optimization}
Given intent $a$ and audience state $b$, the speaker's problem is:
\[
  \max_{s \in \IS}\; u_S(s) = \max_{s \in \IS}\;
  \bigl[\rs(s, a) + \rs(b, a) + \emerge(s, b, a)\bigr].
\]
Since $\rs(b, a)$ is fixed (ethos is given), this reduces to:
\[
  \max_{s \in \IS}\; \bigl[\rs(s, a) + \emerge(s, b, a)\bigr].
\]
\end{definition}

\begin{theorem}[Ethos--Pathos Tradeoff]
\label{thm:ethos-pathos-tradeoff}
If the speaker can choose between a high-logos signal $s_L$ (with
$\rs(s_L, a)$ large but $\emerge(s_L, b, a)$ small or negative) and a
high-pathos signal $s_P$ (with $\rs(s_P, a)$ small but $\emerge(s_P, b, a)$
large), the optimal choice depends on the audience:
\begin{itemize}
  \item \textbf{Friendly audience} ($\rs(b, a) \gg 0$): prefer $s_L$.
        The audience already agrees; reinforce with logic.
  \item \textbf{Hostile audience} ($\rs(b, a) \ll 0$): prefer $s_P$.
        Logic will be rejected; create new meaning through emergence.
  \item \textbf{Neutral audience} ($\rs(b, a) \approx 0$): balance logos
        and pathos.
\end{itemize}
\end{theorem}

\begin{proof}
For a friendly audience, ethos is large and positive; adding logos increases
the total payoff linearly.  For a hostile audience, ethos is large and
negative; even large logos may not overcome the deficit, so the speaker must
rely on emergence (pathos) to create positive resonance that bypasses the
audience's hostility.
\end{proof}

\begin{example}[Martin Luther King Jr.'s ``I Have a Dream'']
\label{ex:mlk}
King faced a mixed audience: supporters (high ethos), opponents (low or
negative ethos), and a neutral national television audience.  His strategy
was \emph{pathos-dominant}: the repeated phrase ``I have a dream'' composed
with the audience's existing knowledge of American ideals to create massive
emergence.  The logos was relatively simple (equality is good, segregation is
bad); the ethos was mixed.  But the pathos---the emergent meaning of the
American dream recomposed with the reality of segregation---was overwhelming.

In IDT terms: $\emerge(s_{\text{dream}}, b_{\text{audience}}, a_{\text{equality}})
\gg |\rs(b_{\text{audience}}, a_{\text{equality}})|$.  Pathos dominated both
logos and ethos.
\end{example}


\section{Rhetoric in Networks: Viral Persuasion}

When rhetoric operates not in a dyad (speaker--audience) but in a
\emph{network}, the dynamics change qualitatively.  A signal does not just
compose with one audience; it cascades through a population, composing with
each host in turn.

\begin{definition}[Viral Spread]
\label{def:viral-spread}
The \emph{viral spread} of idea $v$ through host $a$ is
$\text{Viral}(v, a) := v \compose a$.  The result is the ``mutated'' form
of the idea after passing through host $a$.

\begin{lstlisting}[caption={Lean 4: Viral spread}]
def viralSpread (v a : I) : I := compose v a
\end{lstlisting}
\end{definition}

\begin{theorem}[Viral Enrichment]
\label{thm:viral-enrichment}
Viral spread always enriches: $w(\text{Viral}(v, a)) \geq w(v)$.
The viral idea can only gain weight through transmission.

\begin{lstlisting}[caption={Lean 4: Viral enriches}]
theorem viral_enriches (v a : I) :
    rs (viralSpread v a) (viralSpread v a) ≥ rs v v :=
  compose_enriches' v a
\end{lstlisting}
\end{theorem}

\begin{proof}
$w(v \compose a)^2 \geq w(v)^2$ by E4.  This is just the weight ratchet
applied to viral transmission.

In plain language: a viral idea gets \emph{heavier} with each host it passes
through, because each host's cognitive substance is composed onto it.  This
is why viral ideas become more elaborate and emotionally charged as they
spread: each retransmission adds the retransmitter's substance.  A simple
observation becomes a complex narrative with emotional overtones after
passing through many hosts---not because anyone deliberately made it more
complex, but because composition is enriching by axiom.
\end{proof}

\begin{theorem}[Viral Mutation]
\label{thm:viral-mutation}
The ``mutation'' introduced by host $a$ during viral transmission is:
\[
  \text{Mut}(v, a, c) := \rs(v \compose a, c) - \rs(v, c) = \rs(a, c) +
  \emerge(v, a, c).
\]
The mutation has a predictable component ($\rs(a, c)$---the host's direct
contribution) and an unpredictable component ($\emerge(v, a, c)$---the
emergent distortion).

\begin{lstlisting}[caption={Lean 4: Viral mutation}]
theorem viral_mutation (v a c : I) :
    rs (viralSpread v a) c - rs v c =
    rs a c + emergence v a c := by
  unfold viralSpread; linarith [rs_compose_eq v a c]
\end{lstlisting}
\end{theorem}

\begin{proof}
By the emergence decomposition: $\rs(v \compose a, c) = \rs(v, c) + \rs(a, c)
+ \emerge(v, a, c)$.  Subtracting $\rs(v, c)$ gives the mutation.

The mutation theorem explains why ``telephone game'' dynamics occur in
information networks: each host introduces both their own perspective (the
$\rs(a, c)$ term) and unpredictable emergent distortion (the $\emerge(v, a, c)$
term).  After many hosts, the accumulated mutations can completely transform
the original message.  The weight ratchet ensures that the transformed message
is \emph{heavier} than the original, but it may resonate with very different
ideas.
\end{proof}

\begin{theorem}[Perfect Transmission]
\label{thm:perfect-transmission}
If $\emerge(v, a, c) = 0$ for all $c$, then the viral mutation reduces to
the host's direct contribution: $\text{Mut}(v, a, c) = \rs(a, c)$.  The
message is transmitted without emergent distortion.

\begin{lstlisting}[caption={Lean 4: Perfect transmission}]
theorem perfect_transmission (v a c : I) (h : emergence v a c = 0) :
    rs (viralSpread v a) c - rs v c = rs a c := by
  unfold viralSpread; linarith [rs_compose_eq v a c]
\end{lstlisting}
\end{theorem}

\begin{proof}
Set $\emerge(v, a, c) = 0$ in the mutation formula.  Perfect transmission
occurs when the composition of the viral idea with the host produces no
emergent meaning---the host faithfully relays the message, adding only their
own direct resonance.

In practice, perfect transmission is rare: most compositions produce nonzero
emergence.  This is the formal explanation for the ubiquity of information
distortion in networks.  It is also the reason why institutions invest
heavily in communication protocols, style guides, and standardized
vocabularies: these are attempts to minimize the emergence term during
transmission, keeping the message as close to the original as possible.
\end{proof}

\begin{interpretation}[Rhetoric as Network Strategy]
\label{interp:rhetoric-network}
The viral spread framework connects rhetoric to network science.  A
rhetorician designing a message for viral spread must consider not just
the immediate audience but the entire cascade of future hosts.  The
optimal signal is one that:
\begin{enumerate}
  \item Has high logos (resonates with the intent: $\rs(s, a)$ large).
  \item Has high identification with the first audience (Burke's
        condition: $\rs(s, b)$ large).
  \item Produces positive emergence at each cascade step (the mutation
        is ``beneficial'' at each host: $\emerge(s, b_k, a) > 0$ for
        all hosts $b_k$ in the cascade).
  \item Is \emph{robust} to mutation: the emergent distortions introduced
        by diverse hosts do not accumulate destructively.
\end{enumerate}

Condition (4) is the hardest to satisfy and explains why most viral content
loses its original meaning rapidly.  Only messages with very specific
algebraic properties---those whose emergence terms are consistently positive
across diverse hosts---survive the cascade intact.  Religious texts,
memorable slogans, and catchy melodies all share this property: they compose
productively with a wide range of cognitive states, producing positive
emergence at each step.
\end{interpretation}


% ================================================================
% CHAPTER 14: NARRATIVE AS COMPOSED MEANING
% ================================================================
\chapter{Narrative as Composed Meaning}
\label{ch:narrative-deep}

\epigraph{``The universe is made of stories, not of atoms.''}%
{Muriel Rukeyser}


\section{Stories as Ordered Compositions}

\begin{definition}[Narrative]
\label{def:narrative}
A \emph{narrative} is an ordered sequence of ideas
$\mathbf{a} = [a_1, a_2, \ldots, a_n]$.  The \emph{narrative composite}
(or \emph{story}) is the composed idea:
\[
  S(\mathbf{a}) := a_1 \compose a_2 \compose \cdots \compose a_n.
\]
The \emph{partial narrative} at step $k$ is:
\[
  S_k := a_1 \compose a_2 \compose \cdots \compose a_k.
\]
By convention, $S_0 := \void$.
\end{definition}

\begin{theorem}[Narrative Weight Growth]
\label{thm:narrative-weight-growth}
The weight of the partial narrative is non-decreasing:
\[
  w(S_0) \leq w(S_1) \leq w(S_2) \leq \cdots \leq w(S_n) = w(S(\mathbf{a})).
\]
\end{theorem}

\begin{proof}
$w(S_{k+1}) = w(S_k \compose a_{k+1}) \geq w(S_k)$ by E4.  Induction
from $k = 0$ to $k = n-1$.

\begin{lstlisting}[caption={Lean 4: Narrative weight growth}]
theorem narrative_weight_growth {S : IdeaticSpace I}
    (elements : Fin n → I) (k : Fin n) :
    rs (partialNarrative elements (k + 1))
       (partialNarrative elements (k + 1)) ≥
    rs (partialNarrative elements k)
       (partialNarrative elements k) := by
  simp [partialNarrative]
  exact compose_enriches' _ _
\end{lstlisting}
\end{proof}

\begin{interpretation}[You Cannot Un-Tell a Story]
\label{interp:untell}
Every scene, every sentence, every word adds weight to the narrative.  A
poorly constructed scene damages the story irrecoverably: subsequent scenes
compose on top of it, but they cannot remove its weight.  This is the
mathematical basis of Chekhov's principle: ``If in the first act you have
hung a pistol on the wall, then in the following one it should be fired.''
Unused narrative elements add weight without contributing to the narrative
arc---dead weight that the story must carry to the end.
\end{interpretation}

\begin{remark}[Bakhtin's Dialogism and Narrative Composition]
\label{rem:bakhtin}
Mikhail Bakhtin (1981) argued that narrative is fundamentally
\emph{dialogic}---every utterance is a response to previous utterances and
an anticipation of future responses.  A novel is not a monologue by the
author but a \emph{polyphony} of voices, each with its own worldview.

The IDT framework gives Bakhtin's dialogism a precise formal structure.
Each narrative element $a_k$ is a ``voice'' with its own weight $w(a_k)$
and resonance profile $\rs(a_k, \cdot)$.  The narrative composite
$S(\mathbf{a}) = a_1 \compose \cdots \compose a_n$ is the composed
polyphony---the result of all voices speaking in sequence.

Bakhtin's key insight was that dialogic interaction creates meaning that
no single voice possesses.  In IDT, this is \emph{emergence}: the plot
value $\sum E_k(d)$ measures the genuinely dialogic meaning---the surplus
created by the \emph{interaction} of voices, beyond what each voice
contributes individually.

The non-commutativity of composition ($a_k \compose a_{k+1} \neq
a_{k+1} \compose a_k$) is the formal expression of Bakhtin's observation
that dialogue is \emph{asymmetric}: who speaks first shapes the context
for who speaks next.  In a Bakhtinian novel, the ordering of voices is not
arbitrary---it is a compositional strategy that maximizes the dialogic
emergence.

\emph{Heteroglossia}---Bakhtin's term for the multiplicity of social
languages within a text---corresponds to high \emph{diversity} among
narrative elements: $\rs(a_i, a_j)$ varies widely across different pairs,
reflecting the different registers, dialects, and worldviews represented in
the text.  By the weight divergence theorem
(Theorem~\ref{thm:weight-divergence}), high diversity produces faster weight
growth: heteroglossic narratives are richer (heavier) than monoglossic ones.
\end{remark}

\begin{remark}[Brooks's ``Reading for the Plot'' and Narrative Desire]
\label{rem:brooks}
Peter Brooks (1984) proposed that narrative is driven by ``desire''---the
reader's drive toward the ending, shaped by textual energies that push
forward (anticipation) and pull backward (retrospection).  In IDT, narrative
desire has a precise correlate: the \emph{emergence gradient}.

Define the \emph{forward emergence} at step $k$ as:
\[
  E_k^+(d) := \emerge(S_{k-1}, a_k, d) \qquad \text{(what the next element adds)}.
\]
Brooks's ``narrative desire'' corresponds to the reader's expectation that
$E_k^+(d)$ will be large---that the next element will create substantial new
meaning.  When emergence is large and positive, the reader experiences
satisfaction (the story is ``going somewhere'').  When emergence is small or
negative, the reader experiences frustration (the story is ``stalling'' or
``going wrong'').

Brooks's concept of the ``dilatory space''---the middle of a narrative where
resolution is deferred---corresponds to a region of moderate emergence:
enough to maintain engagement, but not enough to resolve the central tension.
The formal criterion is:
\[
  0 < E_k^+(d) < E_{\text{climax}}^+(d) \qquad \text{for } k \text{ in the
  dilatory space}.
\]

The climax is the point of maximum emergence: $E_k^+(d) = \max_j E_j^+(d)$.
The dénouement is the region of declining emergence as the narrative resolves:
$E_k^+(d) \to 0$ as $k \to n$.

Brooks's insight that narrative operates through a ``logic of anticipation
and retrospection'' is captured by the cocycle constraint
(Theorem~\ref{thm:narrative-cocycle}): the emergence at step $k$ is
algebraically constrained by the emergence at neighboring steps.  The reader's
anticipation is an implicit model of the cocycle---an intuition about what
kinds of emergence are ``possible'' given the narrative so far.  When the
actual emergence violates these expectations (surprise), the reader
experiences either delight (if the surprise is ``earned''---consistent with
deeper cocycle constraints) or dissatisfaction (if it is ``unearned''---
violating the reader's sense of narrative coherence).
\end{remark}


\section{Plot as Emergence Pattern}

\begin{definition}[Narrative Emergence]
\label{def:narrative-emergence}
The \emph{narrative emergence at step $k$}, observed by $d \in \IS$, is:
\[
  E_k(d) := \emerge(S_{k-1}, a_k, d).
\]
This captures how much genuinely new meaning the $k$-th element creates,
given everything that preceded it, as perceived by observer $d$.
\end{definition}

\begin{definition}[Plot]
\label{def:plot}
The \emph{plot} of narrative $\mathbf{a}$ (relative to observer $d$) is the
sequence of emergence values:
\[
  \text{Plot}(\mathbf{a}, d) := \bigl(E_1(d),\, E_2(d),\, \ldots,\, E_n(d)\bigr).
\]
The plot is a real-valued sequence: it records the rise and fall of emergent
meaning throughout the story.
\end{definition}

\begin{theorem}[Plot Determines Resonance Structure]
\label{thm:plot-resonance}
The observer's total resonance with the narrative composite is:
\[
  \rs(S(\mathbf{a}), d) = \sum_{k=1}^{n} \rs(a_k, d)
  + \sum_{k=2}^{n} E_k(d).
\]
The first sum is the ``literal'' content (direct resonance of each element
with the observer).  The second sum is the \emph{plot value}---the total
emergent meaning.
\end{theorem}

\begin{proof}
By telescoping the emergence decomposition.  For $k \geq 2$:
\[
  \rs(S_k, d) = \rs(S_{k-1} \compose a_k, d)
  = \rs(S_{k-1}, d) + \rs(a_k, d) + E_k(d).
\]
Summing from $k = 1$ to $n$ and using $S_0 = \void$ (so $\rs(S_0, d) = 0$):
\[
  \rs(S_n, d) = \sum_{k=1}^n \rs(a_k, d) + \sum_{k=2}^n E_k(d). \qedhere
\]
\end{proof}

\begin{interpretation}[Plot is What Makes a Story More Than Its Parts]
\label{interp:plot}
The plot value $\sum E_k(d)$ measures the total emergent meaning---the
``extra'' resonance created by the \emph{order} and \emph{interaction} of the
narrative elements, beyond their individual contributions.  A story with zero
plot value ($E_k = 0$ for all $k$) is a mere list: each element stands alone,
and the whole is exactly the sum of its parts.  A story with large positive
plot value achieves a whole greater than the sum---the hallmark of great
narrative.

Negative emergence ($E_k < 0$) corresponds to narrative elements that
\emph{detract} from the story when placed in their particular context.  A
joke at a funeral has negative emergence: its resonance with the listener is
diminished by its context, not enhanced.
\end{interpretation}

\begin{remark}[Natural Language Proof of the Plot--Resonance Theorem]
\label{rem:plot-resonance-intuition}
The plot determines resonance structure theorem
(Theorem~\ref{thm:plot-resonance}) has a beautiful telescoping proof that
reveals the compositional nature of narrative meaning.

The core idea is to ``unpack'' the final composite $S_n$ step by step.
At each step, the emergence decomposition (E3) breaks the resonance with
observer $d$ into three pieces: the resonance of the existing partial
narrative, the resonance of the new element, and the emergence.

\textbf{Step 1.}  $\rs(S_1, d) = \rs(a_1, d)$ (the first element contributes
its resonance directly; $S_0 = \void$ contributes nothing).

\textbf{Step 2.}  $\rs(S_2, d) = \rs(S_1 \compose a_2, d) = \rs(S_1, d) +
\rs(a_2, d) + E_2(d) = \rs(a_1, d) + \rs(a_2, d) + E_2(d)$.

\textbf{Step $k$.}  $\rs(S_k, d) = \rs(S_{k-1}, d) + \rs(a_k, d) + E_k(d)$.

Summing the telescoping series from $k = 1$ to $n$:
\[
  \rs(S_n, d) = \sum_{k=1}^n \rs(a_k, d) + \sum_{k=2}^n E_k(d).
\]

The first sum is the ``literal'' content of the narrative: what each element
contributes on its own, independently of context.  The second sum is the
\emph{plot}---the total emergent meaning created by the specific order in
which elements are composed.

This decomposition explains why a narrative is more than a collection of
facts: the ``facts'' (individual element resonances) account for one part
of the total meaning, and the ``plot'' (accumulated emergence) accounts for
the rest.  A well-constructed plot can double or triple the total resonance
compared to a mere enumeration of the same elements.
\end{remark}


\section{Dramatic Tension}

\begin{definition}[Dramatic Tension]
\label{def:dramatic-tension}
The \emph{dramatic tension} between consecutive elements $a_k$ and $a_{k+1}$
(relative to observer $d$) is:
\[
  T_k(d) := |E_{k+1}(d)| = |\emerge(S_k, a_{k+1}, d)|.
\]
High tension = large emergence (positive or negative) between consecutive
elements.  Low tension = small emergence.
\end{definition}

\begin{remark}[Why Absolute Value?]
Dramatic tension is about the \emph{magnitude} of surprise, not its sign.
A shocking revelation ($E > 0$) and a devastating reversal ($E < 0$) both
create high tension.  A predictable next step ($E \approx 0$) creates low
tension.
\end{remark}

\begin{definition}[The Dramatic Arc]
\label{def:dramatic-arc}
The \emph{dramatic arc} of a narrative is the sequence of tensions:
\[
  \text{Arc}(\mathbf{a}, d) := (T_1(d), T_2(d), \ldots, T_{n-1}(d)).
\]
Classical narrative structure (Freytag's pyramid) prescribes:
\begin{enumerate}
  \item \textbf{Exposition}: low tension ($T_k$ small).
  \item \textbf{Rising action}: increasing tension ($T_{k+1} > T_k$).
  \item \textbf{Climax}: maximum tension ($T_k$ at peak).
  \item \textbf{Falling action}: decreasing tension ($T_{k+1} < T_k$).
  \item \textbf{Dénouement}: low tension ($T_k$ returns to small values).
\end{enumerate}
\end{definition}

\begin{theorem}[Tension--Weight Relation]
\label{thm:tension-weight}
High dramatic tension implies large weight growth:
\[
  T_k(S_k) = |\emerge(S_k, a_{k+1}, S_k)|
  \leq w(S_{k+1})^2 - w(S_k)^2.
\]
When $E_{k+1}(S_k) > 0$, the tension is bounded by the weight increment:
dramatic tension contributes to narrative weight.
\end{theorem}

\begin{proof}
By the emergence decomposition with $d = S_k$:
\begin{align*}
  w(S_{k+1})^2 &= \rs(S_k \compose a_{k+1}, S_k \compose a_{k+1}) \\
  &\geq \rs(S_k, S_k) + \rs(a_{k+1}, S_k) + \emerge(S_k, a_{k+1}, S_k) \\
  &= w(S_k)^2 + \rs(a_{k+1}, S_k) + E_{k+1}(S_k).
\end{align*}
Since $\rs(a_{k+1}, S_k) \geq 0$ (when $a_{k+1}$ resonates positively with
the existing narrative), we get
$w(S_{k+1})^2 - w(S_k)^2 \geq E_{k+1}(S_k) = T_k(S_k)$ when $E > 0$.
\end{proof}

\begin{remark}[Natural Language Proof of the Tension--Weight Relation]
\label{rem:tension-weight-intuition}
The tension--weight relation reveals a deep connection between aesthetics
and structure: \textbf{dramatic tension costs weight}.  More precisely,
high tension (large emergence between consecutive elements) is bounded by
the weight increment: the story must ``pay'' for dramatic tension with
increased narrative weight.

This has a practical interpretation for writers.  A scene with high dramatic
tension (a surprising revelation, a shocking twist) necessarily adds
substantial weight to the narrative.  This weight must be carried through
all subsequent scenes.  A story that deploys too many high-tension scenes
early on becomes ``top-heavy''---burdened with so much accumulated weight
that later scenes must compose on top of an enormously heavy context,
making it difficult to create additional tension.

The best narrative structures deploy tension \emph{economically}: each
high-tension scene creates enough weight to support the next level of
dramatic structure, but not so much that the narrative becomes unwieldy.
This is the mathematical content of the ``pacing'' advice given in creative
writing workshops: don't put all the big revelations in the first act,
because the weight they create will crush the rest of the story.

Conversely, a story with no tension ($E_k \approx 0$ for all $k$) has
minimal weight growth: each element adds only its own weight, with no
emergence.  Such a story is ``flat''---a mere enumeration of facts.  The
optimal tension profile is therefore a balance: enough emergence to create
meaningful weight growth (the story ``goes somewhere''), but not so much
that the weight becomes unmanageable (the story doesn't collapse under its
own complexity).
\end{remark}

\begin{example}[The Dramatic Tension of \emph{Hamlet}]
\label{ex:hamlet}
Consider the narrative sequence of \emph{Hamlet} in IDT terms:
\begin{itemize}
  \item $a_1$ = ``King is dead'' (exposition).
        $T_0 = |\emerge(\void, a_1, d)|$ is modest (setting the scene).
  \item $a_2$ = ``Ghost appears'' (inciting incident).
        $T_1 = |\emerge(a_1, a_2, d)|$ is large: the emergence of a ghost
        in the context of a dead king creates massive surprise.
  \item $a_3$ = ``Ghost reveals murder'' (rising action).
        $T_2 = |\emerge(a_1 \compose a_2, a_3, d)|$ is very large: the
        emergence of murder accusation, composed with ghost-of-dead-king,
        creates explosive new meaning.
  \item $a_k$ = ``To be or not to be'' (philosophical climax).
        $T_{k-1}$ reaches its peak: the emergence of existential doubt,
        composed with the weight of accumulated treachery and indecision,
        creates the maximum tension of the play.
\end{itemize}
Shakespeare's genius lies in choosing elements $a_k$ that produce maximal
emergence with the existing narrative context $S_{k-1}$---elements that are
\emph{surprising yet fitting}, to use Aristotle's criterion.
\end{example}

\begin{example}[The Narrative Arc of Scientific Discovery]
\label{ex:scientific-narrative}
Scientific papers follow a narrative structure that can be analyzed through
IDT:
\begin{itemize}
  \item $a_1$ = Introduction (background, context).
        $E_1(d) \approx 0$: the reader is oriented but not surprised.
  \item $a_2$ = The Problem (what is unknown).
        $E_2(d) > 0$: emergence of the gap between what is known and what is
        not, creating intellectual tension.
  \item $a_3$ = Methods (how the problem is attacked).
        $E_3(d)$ moderate: the approach may be novel (high emergence) or
        standard (low emergence).
  \item $a_4$ = Results (what was found).
        $E_4(d)$ peaks: the emergence of empirical findings composed with the
        methods and the problem creates the maximal new meaning.
  \item $a_5$ = Discussion (what it means).
        $E_5(d)$ declines: the results are interpreted, reducing tension by
        connecting findings to existing knowledge.
\end{itemize}
This is Freytag's pyramid applied to scientific writing: the results section
is the climax, the discussion is the falling action.  The IMRaD format is an
implicit optimization of emergence density across sections.
\end{example}

\begin{remark}[Narrative as Strategic Composition: A Synthesis]
\label{rem:narrative-synthesis}
Chapters~\ref{ch:rhetoric-deep} and \ref{ch:narrative-deep} reveal that
rhetoric and narrative are two aspects of the same mathematical structure:
\emph{strategic composition}.

In rhetoric, the speaker chooses a single signal $s$ to maximize
$\rs(s \compose b, a)$---the payoff from composing $s$ with the audience's
state $b$ relative to the intent $a$.  This is a \emph{one-shot} optimization
over the ideatic space.

In narrative, the author chooses a \emph{sequence} of elements $[a_1, \ldots,
a_n]$ to maximize the total emergence $\sum E_k(d)$---the plot value relative
to the reader $d$.  This is a \emph{sequential} optimization over permutations
of the element set.

The connection is that each narrative step is a rhetorical act: the author
is ``speaking'' element $a_k$ to a ``listener'' whose state is the partial
narrative $S_{k-1}$.  The emergence $E_k(d) = \emerge(S_{k-1}, a_k, d)$ is
the ``pathos'' of the $k$-th rhetorical act.  Narrative is iterated rhetoric.

This perspective unifies classical rhetoric (Aristotle, Quintilian) with
classical narratology (Propp, Barthes, Genette) within a single mathematical
framework.  The formal theorems---weight ratchet, cocycle constraint,
NP-hardness of optimal ordering---apply simultaneously to both domains
because both are instances of strategic composition in the ideatic monoid.
\end{remark}


\section{Narrative Order Matters}

\begin{theorem}[Narrative Non-Commutativity]
\label{thm:narrative-noncommutative}
Reordering the elements of a narrative changes its composite:
\[
  [a_1, a_2, \ldots, a_n] \neq [a_{\sigma(1)}, a_{\sigma(2)}, \ldots, a_{\sigma(n)}]
\]
in general (for permutations $\sigma \neq \text{id}$), and the plot and
dramatic arc change accordingly.
\end{theorem}

\begin{proof}
$S(\mathbf{a}) = a_1 \compose \cdots \compose a_n \neq
a_{\sigma(1)} \compose \cdots \compose a_{\sigma(n)} = S(\sigma \cdot \mathbf{a})$
when composition is non-commutative.  The emergence terms change because
$\emerge(S_{k-1}, a_k, d)$ depends on the partial narrative $S_{k-1}$, which
is different under different orderings.
\end{proof}

\begin{interpretation}[Why Flashbacks Work]
\label{interp:flashbacks}
A flashback reorders the narrative: an event that occurred early in the
fabula (story-world chronology) is placed late in the sjuzhet (narrative
presentation order).  By Theorem~\ref{thm:narrative-noncommutative}, this
changes the composite and the entire emergence pattern.

A flashback ``works'' when $\emerge(S_{k-1}, a_{\text{flashback}}, d)$ is
large---when the revealed past event, composed with everything the reader
now knows, creates massive emergence.  The same event presented in
chronological order would have small emergence (because the context $S_{k-1}$
would be thinner).  The flashback \emph{engineers} large emergence by
controlling the composition order.

This is agenda-setting (Chapter~\ref{ch:voting-deep}) applied to narrative:
the author controls the order of ``idea composition'' to maximize dramatic
effect.
\end{interpretation}

\begin{remark}[Genette's Narrative Time and IDT Composition Order]
\label{rem:genette}
Gérard Genette (1980) distinguished three dimensions of narrative time:
\emph{order} (the sequence of events as presented vs.\ as occurred),
\emph{duration} (how much narrative time is devoted to each event), and
\emph{frequency} (how often an event is narrated).

In IDT:
\begin{itemize}
  \item \textbf{Order} is the composition permutation $\sigma$.  Different
        orderings produce different narrative composites
        (Theorem~\ref{thm:narrative-noncommutative}) and different emergence
        patterns.  \emph{Analepsis} (flashback) and \emph{prolepsis}
        (flash-forward) are specific permutations chosen to maximize
        emergence at the point of insertion.

  \item \textbf{Duration} corresponds to the \emph{weight} of each narrative
        element $w(a_k)$.  A scene described in great detail has high weight;
        a summary has low weight.  By the weight ratchet, a heavy scene
        contributes more to the narrative composite than a light summary.
        This is why ``showing'' is more powerful than ``telling'': showing
        produces heavier elements that contribute more weight to the composite.

  \item \textbf{Frequency} corresponds to \emph{iterated composition}.
        An event narrated twice ($a_k$ composed at two different points)
        creates two emergence events: $E_k(d)$ at the first telling and
        $E_{k'}(d)$ at the second.  By the weight ratchet, the second
        telling composes with a heavier context (including the weight of
        the first telling), potentially producing very different emergence.
        This is why repetition in narrative is not mere redundancy: each
        repetition occurs in a different compositional context, creating
        different emergent meaning.
\end{itemize}
\end{remark}

\begin{theorem}[Narrative Reordering Bound]
\label{thm:reordering-bound}
For any two permutations $\pi, \sigma$ of a narrative $\mathbf{a}$, the
difference in observer resonance is bounded:
\[
  \bigl|\rs(S_\pi(\mathbf{a}), d) - \rs(S_\sigma(\mathbf{a}), d)\bigr|
  \leq \sum_{k=2}^{n} \bigl|E_k^\pi(d) - E_k^\sigma(d)\bigr|,
\]
where $E_k^\pi$ and $E_k^\sigma$ are the emergence sequences under the
two orderings.  The difference is entirely in the \emph{plot}---the literal
content $\sum \rs(a_k, d)$ is the same for all permutations.
\end{theorem}

\begin{proof}
By Theorem~\ref{thm:plot-resonance}, $\rs(S_\pi, d) = \sum_k \rs(a_k, d)
+ \sum_k E_k^\pi(d)$ and $\rs(S_\sigma, d) = \sum_k \rs(a_k, d) +
\sum_k E_k^\sigma(d)$.  The literal content sums are identical (they sum
over the same elements in different order, and addition is commutative).
Subtracting: $|\rs(S_\pi, d) - \rs(S_\sigma, d)| = |\sum_k (E_k^\pi(d) -
E_k^\sigma(d))| \leq \sum_k |E_k^\pi(d) - E_k^\sigma(d)|$.

In plain language: two tellings of the same story (same events, different
order) differ only in their \emph{plot value}---the emergence pattern.
The ``facts'' are the same; the ``meaning'' (emergent resonance) differs.
This is the mathematical content of the literary-critical truism that
``the same story can be told many ways'': the same narrative elements can
be composed in different orders, producing different emergence patterns
and therefore different total meanings.  But the difference is bounded: it
cannot exceed the sum of emergence differences across all steps.
\end{proof}


\section{Narrative Density and Economy}

\begin{definition}[Narrative Density]
\label{def:narrative-density}
The \emph{narrative density} of story $\mathbf{a}$ is:
\[
  \rho(\mathbf{a}) := \frac{w(S(\mathbf{a}))^2}{n} =
  \frac{\rs(S(\mathbf{a}), S(\mathbf{a}))}{n}.
\]
Weight per element.
\end{definition}

\begin{theorem}[Density Lower Bound]
\label{thm:density-lower-bound}
$\rho(\mathbf{a}) \geq w(a_1)^2 / n$.
\end{theorem}

\begin{proof}
$w(S(\mathbf{a})) \geq w(a_1)$ by repeated application of E4.
\end{proof}

\begin{definition}[Emergence Density]
\label{def:emergence-density}
The \emph{emergence density} is:
\[
  \rho_E(\mathbf{a}, d) := \frac{\sum_{k=2}^n E_k(d)}{n - 1}.
\]
Average emergence per transition.  This measures the ``surprises per scene.''
\end{definition}

\begin{theorem}[Poetry vs.\ Prose: A Density Criterion]
\label{thm:poetry-prose}
If we define \emph{poetry} as narrative with emergence density above threshold
$\rho_0$ and \emph{prose} as narrative with density below $\rho_0$, then:
\begin{enumerate}
  \item Poetry has higher weight per element than prose (of the same length).
  \item Poetry tolerates shorter sequences (fewer elements for the same total
        weight).
  \item Poetry requires higher-emergence transitions, which correlates with
        greater ambiguity and difficulty of interpretation.
\end{enumerate}
\end{theorem}

\begin{proof}
(1): Total weight $= \sum \rs(a_k, d) + \sum E_k(d)$.  Higher $\rho_E$
means more emergence per element, hence more weight per element.

(2): For target weight $W$, the number of elements needed is $n \approx W /
\rho$, which is smaller when $\rho$ (and hence $\rho_E$) is larger.

(3): High emergence means the composed meaning deviates significantly from
the ``literal'' sum of parts, making interpretation harder.
\end{proof}

\begin{interpretation}[Coleridge's ``Best Words in the Best Order'']
\label{interp:coleridge}
Coleridge defined poetry as ``the best words in the best order.''  In IDT,
``best words'' means elements with high individual weight ($w(a_k)$ large)
and ``best order'' means the permutation that maximizes total emergence
($\sum E_k$ maximal).  The poet's craft is the joint optimization of element
selection and element ordering---a problem that is computationally hard
(it involves searching over both the elements and their permutations) but
aesthetically solvable by human intuition.
\end{interpretation}

\begin{remark}[The Density Spectrum: From Journalism to Poetry]
\label{rem:density-spectrum}
The emergence density metric $\rho_E$ provides a formal spectrum from
low-density to high-density narrative:

\begin{itemize}
  \item \textbf{Journalism} ($\rho_E \approx 0$): the goal is to convey
        information with minimal emergence.  Each sentence adds its direct
        resonance ($\rs(a_k, d)$) with minimal interaction effects.  The
        ``inverted pyramid'' structure minimizes emergence by placing the most
        important elements first and the least important last, so that each
        element can be understood independently.

  \item \textbf{Prose fiction} ($\rho_E$ moderate): the plot creates
        emergence between narrative elements, but each element also carries
        substantial direct resonance.  The balance between literal meaning
        and emergent meaning is roughly even.

  \item \textbf{Poetry} ($\rho_E$ high): most of the meaning is emergent.
        Individual words may have low direct resonance with the theme, but
        their \emph{interactions}---the emergence created by their specific
        juxtaposition---carry the bulk of the meaning.  This is why poetry
        is ``untranslatable'': translation preserves the direct resonance
        $\rs(a_k, d)$ of each word but destroys the emergence pattern
        $E_k(d)$, which depends on the specific compositional structure
        of the source language.

  \item \textbf{Music} ($\rho_E$ maximal): in purely instrumental music,
        the ``narrative elements'' (notes, chords, rhythmic patterns) have
        minimal direct resonance with any semantic content, but their
        temporal composition creates massive emergence.  Music is
        ``all plot, no literal content''---the extreme case of emergence-
        dominated narrative.
\end{itemize}
\end{remark}


\section{The Cocycle and Narrative Coherence}

\begin{theorem}[Narrative Cocycle Constraint]
\label{thm:narrative-cocycle}
The cocycle condition constrains the plot: for any three consecutive elements
$a_k, a_{k+1}, a_{k+2}$ and observer $d$:
\[
  \emerge(a_k, a_{k+1}, d) + \emerge(a_k \compose a_{k+1}, a_{k+2}, d) =
  \emerge(a_{k+1}, a_{k+2}, d) + \emerge(a_k, a_{k+1} \compose a_{k+2}, d).
\]
\end{theorem}

\begin{proof}
Direct application of the cocycle condition (E3 consequence) with
$a = a_k$, $b = a_{k+1}$, $c = a_{k+2}$.

\begin{lstlisting}[caption={Lean 4: Narrative cocycle}]
theorem narrative_cocycle {S : IdeaticSpace I}
    (a₁ a₂ a₃ d : I) :
    emergence a₁ a₂ d + emergence (compose a₁ a₂) a₃ d =
    emergence a₂ a₃ d + emergence a₁ (compose a₂ a₃) d :=
  cocycle_condition a₁ a₂ a₃ d
\end{lstlisting}
\end{proof}

\begin{interpretation}[Narrative Coherence as Cocycle Satisfaction]
\label{interp:coherence}
The cocycle constraint is a \emph{coherence condition} on the plot.  It says
that the emergence of ``$a_k$ then $a_{k+1}$, then $a_{k+2}$'' is
algebraically related to the emergence of ``$a_{k+1}$ then $a_{k+2}$'' and
``$a_k$ then ($a_{k+1}$ composed with $a_{k+2}$).''  A narrative that
violates this constraint is internally incoherent---it has plot holes.

In practice, the cocycle is always satisfied (it follows from the axioms),
but the \emph{craft} of narrative lies in choosing elements whose cocycle
relations produce aesthetically pleasing patterns: the emergence of the
climax should be ``prepared'' by the rising action in a way that satisfies
the cocycle.  A deus ex machina is an element $a_k$ whose emergence is large
but whose cocycle relation with the preceding elements is ``unearned''---
the audience senses that the plot structure is algebraically incoherent,
even if they cannot articulate the cocycle condition.
\end{interpretation}


\section{The Narrator's Optimization Problem}

\begin{definition}[Optimal Narrative]
\label{def:optimal-narrative}
Given a set of narrative elements $\{a_1, \ldots, a_n\}$, a target audience
$d$, and a dramatic arc template $\mathbf{T}^* = (T_1^*, \ldots, T_{n-1}^*)$,
the \emph{narrator's optimization problem} is:
\[
  \min_{\sigma \in S_n} \sum_{k=1}^{n-1}
  \bigl|T_k(d; \sigma) - T_k^*\bigr|^2,
\]
where $T_k(d; \sigma)$ is the dramatic tension at step $k$ under
permutation $\sigma$.
\end{definition}

\begin{theorem}[NP-Hardness of Optimal Narrative]
\label{thm:narrative-np-hard}
The narrator's optimization problem is NP-hard in general (by reduction
from the traveling salesman problem on the ``distance'' $|T_k - T_k^*|$
between successive elements).
\end{theorem}

\begin{proof}[Proof sketch]
Construct an ideatic space where $n$ elements form a ``weight graph'' with
edge weights corresponding to emergence magnitudes.  The optimal ordering
that matches a target tension profile maps to a minimum-weight Hamiltonian
path in this graph, which is NP-hard.
\end{proof}

\begin{interpretation}[Why Storytelling Is Hard]
\label{interp:storytelling-hard}
The NP-hardness of optimal narrative explains why great storytelling is rare:
even given the ``right'' elements, finding the right order is computationally
intractable.  Human narrators rely on heuristics, intuition, and aesthetic
sense---approximate algorithms for a problem that has no efficient exact
solution.

This also explains why revision is essential to good writing: the first
draft is an initial permutation, and revision is a local search through
neighboring permutations, each evaluated by the writer's aesthetic sense
(an approximate oracle for the emergence pattern).
\end{interpretation}

\begin{remark}[Greedy Narrative Algorithms]
\label{rem:greedy-narrative}
Although optimal narrative is NP-hard, \emph{greedy} algorithms provide
useful approximations.  A greedy narrator, at each step, chooses the element
$a_k$ from the remaining set that maximizes the immediate emergence
$E_k(d) = \emerge(S_{k-1}, a_k, d)$.

The greedy algorithm has an intuitive interpretation: at each moment in the
story, choose the ``most surprising'' next element given the current context.
This is often good advice for thriller writers but poor advice for literary
novelists, because the greedy algorithm maximizes \emph{local} emergence at
the expense of \emph{global} structure.  A locally surprising sequence may
have poor global coherence (the cocycle constraints are not satisfied in a
way that produces a satisfying overall arc).

More sophisticated narrative algorithms would ``look ahead'' multiple steps,
considering how the choice at step $k$ affects the emergence at steps
$k+1, k+2, \ldots$  This look-ahead is constrained by the cocycle:
$E_k$ and $E_{k+1}$ are algebraically linked, so choosing to maximize
$E_k$ may reduce or increase $E_{k+1}$ depending on the cocycle structure.

The observation that human narrators seem to use a combination of greedy
local choices and global structural templates (genre conventions, the
three-act structure, the hero's journey) suggests that human narrative
intuition operates as a \emph{hybrid algorithm}: greedy at the scene level,
template-matching at the act level, and cocycle-satisfying at the whole-work
level.
\end{remark}


\section{The Composition of Sub-Narratives}

\begin{definition}[Sub-Narrative and Macro-Narrative]
\label{def:sub-narrative}
A \emph{sub-narrative} is a contiguous subsequence $[a_j, a_{j+1}, \ldots, a_k]$.
Its composite is $S_{j,k} := a_j \compose \cdots \compose a_k$.

A \emph{macro-narrative} treats each sub-narrative as a single element:
if the full narrative is divided into chapters $C_1, C_2, \ldots, C_m$
with composites $S(C_1), S(C_2), \ldots, S(C_m)$, then the macro-narrative
is $[S(C_1), S(C_2), \ldots, S(C_m)]$.
\end{definition}

\begin{theorem}[Macro-Narrative Equals Full Narrative]
\label{thm:macro-equals-full}
By associativity:
\[
  S(C_1) \compose S(C_2) \compose \cdots \compose S(C_m)
  = a_1 \compose a_2 \compose \cdots \compose a_n = S(\mathbf{a}).
\]
The macro-narrative composite equals the full narrative composite.
\end{theorem}

\begin{proof}
Immediate from associativity of $\compose$.

\begin{lstlisting}[caption={Lean 4: Macro-narrative associativity}]
theorem macro_narrative_eq_full {S : IdeaticSpace I}
    (ch1 ch2 : I) (rest : I) :
    compose (compose ch1 ch2) rest =
    compose ch1 (compose ch2 rest) := by
  -- associativity of compose
  rfl  -- if compose is defined as monoid operation
\end{lstlisting}
\end{proof}

\begin{theorem}[Chapter-Level Emergence]
\label{thm:chapter-emergence}
The emergence between consecutive chapters is:
\[
  E_{\text{ch}}(C_j, C_{j+1}, d) :=
  \emerge\!\left(S(C_1) \compose \cdots \compose S(C_j),\;
                 S(C_{j+1}),\; d\right).
\]
This ``chapter-level emergence'' captures the macro-dramatic tension:
how much new meaning each chapter creates given everything before it.
\end{theorem}

\begin{interpretation}[Fractal Narrative Structure]
\label{interp:fractal}
Great narratives exhibit emergence at multiple scales simultaneously:
sentence-level emergence (word choice), scene-level emergence (plot beats),
chapter-level emergence (major revelations), and whole-work emergence
(thematic resonance).  The cocycle condition ensures that these levels are
algebraically consistent: the macro-emergence is constrained by the
micro-emergence through the cocycle.

This fractal structure is what distinguishes a masterwork from a competent
plot: in a masterwork, the emergence pattern is rich at every scale of
analysis.  In IDT, this corresponds to high $|E_k(d)|$ at every level of
the sub-narrative decomposition.
\end{interpretation}


\section{Narrative Networks and Intertextuality}

Narratives do not exist in isolation.  Every story composes with the reader's
knowledge of other stories---a phenomenon that literary theory calls
\emph{intertextuality} (Kristeva, 1966).

\begin{definition}[Intertextual Composition]
\label{def:intertextuality}
The \emph{intertextual reading} of narrative $\mathbf{a}$ by a reader whose
literary background includes narratives $\mathbf{b}_1, \ldots, \mathbf{b}_m$
is:
\[
  R_{\text{inter}} := S(\mathbf{a}) \compose S(\mathbf{b}_1) \compose
  \cdots \compose S(\mathbf{b}_m).
\]
The intertextual resonance with theme $d$ is:
\[
  \rs(R_{\text{inter}}, d) = \rs(S(\mathbf{a}), d) +
  \sum_{j=1}^{m} \rs(S(\mathbf{b}_j), d) +
  \text{emergence terms}.
\]
\end{definition}

\begin{theorem}[Intertextual Enrichment]
\label{thm:intertextual-enrichment}
The intertextual reading is always richer than the isolated reading:
\[
  w(R_{\text{inter}}) \geq w(S(\mathbf{a})).
\]
A reader who brings more literary background to a text extracts more meaning
from it.
\end{theorem}

\begin{proof}
By repeated application of E4: each additional narrative composed onto the
reading state only increases weight.

In plain language: a reader who has read Homer, Dante, and Shakespeare
will extract more meaning from Joyce's \emph{Ulysses} than a reader who
has read none of them.  This is because the intertextual references in
\emph{Ulysses} create emergence with the reader's knowledge of the
referenced works.  The more prior works are ``composed'' into the reader's
state, the richer the emergence with new works.

This is the formal justification for the literary canon: canonical works are
those that maximize intertextual emergence---they compose productively with
the largest number of other works, creating the richest possible reading
experience for a well-read reader.
\end{proof}

\begin{interpretation}[The Narrative as Composed Meaning: A Final Synthesis]
\label{interp:final-narrative}
The theory developed in this chapter reveals narrative as the most complex
form of strategic composition.  Where a single rhetorical act composes one
signal with one audience state, a narrative composes $n$ elements in
sequence, producing $n-1$ emergence events that collectively constitute
the \emph{plot}.

The key results are:

\begin{enumerate}
  \item \textbf{Narrative weight grows monotonically}
        (Theorem~\ref{thm:narrative-weight-growth}): stories can only gain
        weight, never lose it.  Every element adds substance.

  \item \textbf{Plot determines resonance structure}
        (Theorem~\ref{thm:plot-resonance}): the total meaning of a story
        decomposes into literal content (element resonances) and plot
        (emergence sums).

  \item \textbf{Narrative order matters}
        (Theorem~\ref{thm:narrative-noncommutative}): the same elements in
        different orders produce different stories with different meanings.

  \item \textbf{The cocycle constrains plot}
        (Theorem~\ref{thm:narrative-cocycle}): not all emergence patterns
        are possible---the algebraic structure of $\IS$ constrains which
        plots are coherent.

  \item \textbf{Optimal narrative is NP-hard}
        (Theorem~\ref{thm:narrative-np-hard}): finding the best order for
        narrative elements is computationally intractable.

  \item \textbf{Macro-narratives compose associatively}
        (Theorem~\ref{thm:macro-equals-full}): chapter-level structure is
        consistent with scene-level structure.

  \item \textbf{Intertextuality enriches}
        (Theorem~\ref{thm:intertextual-enrichment}): literary background
        enhances the reading experience through compositional weight growth.
\end{enumerate}

Together, these results establish narrative as the paradigmatic example of
IDT in action: an ordered composition of ideas, constrained by algebraic
structure, optimized (imperfectly, through human craft) for maximum
emergence, and enriched by the intertextual weight of the reader's
cognitive history.

This is the deepest sense in which ``the universe is made of stories, not
of atoms'' (Rukeyser): stories are \emph{composed meanings}---elements
connected through the compositional structure of the ideatic monoid, creating
emergent resonance that transcends the sum of parts.  The mathematics of
IDT gives this poetic claim the precision of a theorem.
\end{interpretation}


\section{Recapitulation: Strategic Interaction in the Ideatic Space}

The seven chapters of Vol~II have developed the theory of strategic
interaction in ideatic spaces.  Let us collect the main threads:

\begin{enumerate}
  \item \textbf{Repeated meaning games} (Chapter~\ref{ch:repeated-deep})
        formalize ongoing discourse.  The weight ratchet ensures irreversibility;
        the folk theorem guarantees multiple equilibria (conventions); echo
        chambers arise from iterated self-reinforcement; polarization grows
        when communities compose inwardly rather than bridging.

  \item \textbf{Coalitions of ideas} (Chapter~\ref{ch:coalitions-deep})
        formalize collaborative thought.  The cooperation surplus measures
        synergy; the cocycle prevents fair attribution of emergence; the
        Shapley value provides a (necessarily order-dependent) approximation
        to fairness.

  \item \textbf{Voting and agenda-setting} (Chapter~\ref{ch:voting-deep})
        formalize collective decision-making about composition order.
        Arrow's impossibility, enriched by the cocycle, shows that democratic
        composition is impossible in a strong sense.  Habermas's deliberation
        and Mouffe's agonism are both modelable within the framework.

  \item \textbf{The evolution of meaning} (Chapter~\ref{ch:evolution-deep})
        formalizes memetic dynamics.  Ideas as replicators, the Price equation
        for cultural evolution, inclusive fitness, Dollo's irreversibility,
        and the fundamental asymmetry between biological and cultural
        evolution (the Lamarckian ratchet).

  \item \textbf{The Red Queen of meaning} (Chapter~\ref{ch:red-queen-deep})
        formalizes coevolutionary dynamics.  Fitness landscapes shift through
        composition; arms races escalate through the weight ratchet;
        speciation occurs through compositional isolation; convergent
        evolution arises through shared environmental composition.

  \item \textbf{Rhetoric as strategic composition}
        (Chapter~\ref{ch:rhetoric-deep}) formalizes persuasion.  The
        logos--ethos--pathos decomposition is exact (not metaphorical);
        Burke's identification is high cross-resonance; Perelman's universal
        audience is the neutral observer; viral spread creates enrichment
        cascades with emergent mutations.

  \item \textbf{Narrative as composed meaning}
        (Chapter~\ref{ch:narrative-deep}) formalizes storytelling.  Plot is
        emergence; dramatic tension is emergence magnitude; the cocycle
        constrains coherence; optimal narrative is NP-hard; intertextuality
        enriches through compositional weight growth; Bakhtin's dialogism
        is the polyphonic composition of heteroglossic voices.
\end{enumerate}

The unifying theme is that \textbf{strategic interaction in ideatic spaces
is governed by three structural principles}: the weight ratchet (composition
only enriches), the cocycle constraint (emergence terms are algebraically
linked), and non-commutativity (order matters).  These three principles,
formalized in Vol~I's axioms E1--E8 and proved in the Lean formalization,
generate the entire theory of strategic meaning-making developed here.
